% NetHack 5.0  Guidebook.tex $NHDT-Date: 1782274458 2026/06/23 23:14:18 $  $NHDT-Branch: NetHack-5.0 $:$NHDT-Revision: 1.630 $ */
\documentclass[titlepage]{article}
\usepackage[hidelinks]{hyperref}
\usepackage{longtable}
\usepackage{enumitem}
\usepackage[a4paper, text={160mm, 220mm}, centering]{geometry}

\setlist[description]{leftmargin=30mm, topsep=2mm, partopsep=0mm, parsep=0mm, itemsep=1mm, labelwidth=28mm, labelsep=2mm}

\hyphenation{CRASHREPORTURL}

\begin{document}
%
% input file: guidebook.mn
%
%.ds h0 "
%.ds h1 %.ds h2 \%
%.ds f0 "

%.mt
\title{\LARGE A Guide to the Mazes of Menace:\\
\Large Guidebook for \textit{NetHack}}

%.au
\author{Original version - Eric S. Raymond\\
(Edited and expanded for 5.0.0 by Mike Stephenson and others)}
%DO NOT REMOVE NH_DATESUB \date{Date(%B %-d, %Y)}
\date{July 4, 2026}

\maketitle

%.pg
%.hn 1
\section{Introduction}

%.pg

Recently, you have begun to find yourself unfulfilled and distant
in your daily occupation.  Strange dreams of prospecting, stealing,
crusading, and combat have haunted you in your sleep for many months,
but you aren't sure of the reason.  You wonder whether you have in
fact been having those dreams all your life, and somehow managed to
forget about them until now.  Some nights you awaken suddenly
and cry out, terrified at the vivid recollection of the strange and
powerful creatures that seem to be lurking behind every corner of the
dungeon in your dream.  Could these details haunting your dreams be real?
As each night passes, you feel the desire to enter the mysterious caverns
near the ruins grow stronger.  Each morning, however, you quickly put
the idea out of your head as you recall the tales of those who entered
the caverns before you and did not return.  Eventually you can resist
the yearning to seek out the fantastic place in your dreams no longer.
After all, when other adventurers came back this way after spending time
in the caverns, they usually seemed better off than when they passed
through the first time.  And who was to say that all of those who did
not return had not just kept going?
%.pg

Asking around, you hear about a bauble, called the Amulet of Yendor by some,
which, if you can find it, will bring you great wealth.  One legend you were
told even mentioned that the one who finds the amulet will be granted
immortality by the gods.  The amulet is rumored to be somewhere beyond the
Valley of Gehennom, deep within the Mazes of Menace.  Upon hearing the
legends, you immediately realize that there is some profound and
undiscovered reason that you are to descend into the caverns and seek
out that amulet of which they spoke.  Even if the rumors of the amulet's
powers are untrue, you decide that you should at least be able to sell the
tales of your adventures to the local minstrels for a tidy sum, especially
if you encounter any of the terrifying and magical creatures of
your dreams along the way.  You spend one last night fortifying yourself
at the local inn, becoming more and more depressed as you watch the odds
of your success being posted on the inn's walls getting lower and lower.

%.pg
\noindent In the morning you awake, collect your belongings, and
set off for the dungeon.  After several days of uneventful
travel, you see the ancient ruins that mark the entrance to the
Mazes of Menace.  It is late at night, so you make camp at the entrance
and spend the night sleeping under the open skies.  In the morning, you
gather your gear, eat what may be your last meal outside, and enter the
dungeon\ldots

%.hn 1
\section{What is going on here?}

%.pg
You have just begun a game of \textit{NetHack}.  Your goal is to grab as much
treasure as you can, retrieve the Amulet of Yendor, and escape the
Mazes of Menace alive.

%.pg
Your abilities and strengths for dealing with the hazards of adventure
will vary with your background and training:

%.pg
%
\begin{description}
\item[Archeologists]
understand dungeons pretty well; this enables them
to move quickly and sneak up on the local nasties.  They start equipped
with the tools for a proper scientific expedition, and are able to read
ancient languages.
%.pg
%
\item[Barbarians]
are warriors out of the hinterland, hardened to battle.
They begin their quests with naught but uncommon strength, a trusty hauberk,
and a great two-handed sword.
%.pg
%
\item[Cavemen \textrm{\textmd{and}} Cavewomen]
start with exceptional strength, but unfortunately, neolithic weapons.
%.pg
%
\item[Healers]
are wise in medicine and apothecary.  They know the
herbs and simples that can restore vitality, ease pain, anesthetize,
and neutralize
poisons; and with their instruments, they can divine a being's state
of health or sickness.  Their medical practice earns them quite reasonable
amounts of money, with which they enter the dungeon.
%.pg
%
\item[Knights]
are distinguished from the common skirmisher by their
devotion to the ideals of chivalry and by the surpassing excellence of
their armor.
%.pg
%
\item[Monks]
are ascetics, who by rigorous practice of physical and mental
disciplines have become capable of fighting as effectively without weapons
as with.  They wear no armor but make up for it with increased mobility.
%.pg
%
\item[Priests \textrm{\textmd{and}} Priestesses]
are clerics militant, crusaders
advancing the cause of righteousness with arms, armor, and arts
thaumaturgic.  Their ability to commune with deities via prayer
occasionally extricates them from peril, but can also put them in it.
%.pg
%
\item[Rangers]
are most at home in the woods, and some say slightly out
of place in a dungeon.  They are, however, experts in archery as well
as tracking and stealthy movement.
%.pg
%
\item[Rogues]
are agile and stealthy thieves, with knowledge of locks,
traps, and poisons.  Their advantage lies in surprise, which they employ
to great advantage.
%.pg
%
\item[Samurai]
are the elite warriors of feudal Nippon.  They are lightly
armored and quick, and wear the %
\textit{dai-sho}, two swords of the deadliest
keenness.
%.pg
%
\item[Tourists]
start out with lots of gold (suitable for shopping with),
a credit card, lots of food, some maps, and an expensive camera.  Most
monsters don't like being photographed.
%.pg
%
\item[Valkyries]
are hardy warrior women.  Their upbringing in the harsh
Northlands makes them strong, inures them to extremes of cold, and instills
in them stealth and cunning.
%.pg
%
\item[Wizards]
start out with a knowledge of magic, a selection of magical
items, and a particular affinity for dweomercraft.  Although seemingly weak
and easy to overcome at first sight, an experienced Wizard is a deadly foe.
\end{description}

%.pg
You may also choose the race of your character (within limits; most
roles have restrictions on which races are eligible for them):

%.pg
%
\begin{description}
\item[Dwarves]
are smaller than humans or elves, but are stocky and solid
individuals.  Dwarves' most notable trait is their great expertise in mining
and metalwork.  Dwarvish armor is said to be second in quality not even to the
mithril armor of the Elves.
%.pg
%
\item[Elves]
are agile, quick, and perceptive; very little of what goes
on will escape an Elf.  The quality of Elven craftsmanship often gives
them an advantage in arms and armor.
%.pg
%
\item[Gnomes]
are smaller than but generally similar to dwarves.  Gnomes are
known to be expert miners, and it is known that a secret underground mine
complex built by this race exists within the Mazes of Menace, filled with
both riches and danger.
%.pg
%
\item[Humans]
are by far the most common race of the surface world, and
are thus the norm to which other races are often compared.  Although
they have no special abilities, they can succeed in any role.
%.pg
%
\item[Orcs]
are a cruel and barbaric race that hate every living thing
(including other orcs).  Above all others, Orcs hate Elves with a passion
unequalled, and will go out of their way to kill one at any opportunity.
The armor and weapons fashioned by the Orcs are typically of inferior quality.
\end{description}

%.hn 1
\section{What do all those things on the screen mean?}
%.pg
On the screen is kept a map of where you have been and what you have
seen on the current dungeon level; as you explore more of the level,
it appears on the screen in front of you.

%.pg
When \textit{NetHack}'s ancestor \textit{rogue} first appeared, its screen
orientation was almost unique among computer fantasy games.  Since
then, screen orientation has become the norm rather than the
exception; \textit{NetHack} continues this fine tradition.  Unlike text
adventure games that accept commands in pseudo-English sentences and
explain the results in words, \textit{NetHack} commands are all one or two
keystrokes and the results are displayed graphically on the screen.  A
minimum screen size of 24 lines by 80 columns is recommended; if the
screen is larger, only a $21\times80$ section will be used for the map.

%.pg
\textit{NetHack} can even be played by blind players, with the assistance of
Braille readers or speech synthesisers.  Instructions for configuring
\textit{NetHack} for the blind are included later in this document.

%.pg
\textit{NetHack} generates a new dungeon every time you play it; even the
authors still find it an entertaining and exciting game despite
having won several times.

%.pg
\textit{NetHack} offers a variety of display options.  The options available to
you will vary from port to port, depending on the capabilities of your
hardware and software, and whether various compile-time options were
enabled when your executable was created.  The three possible display
options are: a monochrome character interface, a color character interface,
and a graphical interface using small pictures called tiles.  The two
character interfaces allow fonts with other characters to be substituted,
but the default assignments use standard ASCII characters to represent
everything.  There is no difference between the various display options
with respect to game play.  Because we cannot reproduce the tiles or
colors in the Guidebook, and because it is common to all ports, we will
use the default ASCII characters from the monochrome character display
when referring to things you might see on the screen during your game.
%.pg
In order to understand what is going on in \textit{NetHack}, first you must
understand what \textit{NetHack} is doing with the screen.  The \textit{NetHack}
screen replaces the ``You see \ldots'' descriptions of text adventure games.
Figure 1 is a sample of what a \textit{NetHack} screen might look like.
The way the screen looks for you depends on your platform.
%.BR 2

% (Either generated by hand or else the composite of two different
% situations.  Originally the character had only reached a second room
% (unchanged here) by turn 257 (now changed to 752) and was already
% Weak from hunger (now changed to just Hungry) and also lacked any of
% Tourist's starting gold.  Confusion is added to include a condition.)
%
% Width is forced to match similar figure in Guidebook.mn where it is
% constrained by the margins of plain text output (Guidebook.txt).
\vbox{
\begin{verbatim}
        The bat bites!

                ------
                |....|    ----------
                |.<..|####...@...$.|
                |....-#   |...B....+
                |....|    |.d......|
                ------    -------|--



        Player the Rambler   St:12 Dx:7 Co:18 In:11 Wi:9 Ch:15 Neutral
        Dlvl:1 $:993 HP:9(12) Pw:3(3) AC:10 Exp:1/19 T:752 Hungry Conf
\end{verbatim}
\begin{center}
Figure 1
\end{center}
}

% 3-line status includes trailing spaces to force the width to match the
% 2-line data above; unlike Guidebook.pm, we can't add a trailing comment
% to make them visible
\vbox{
\begin{verbatim}
        Player the Rambler   St:12 Dx:7 Co:18 In:11 Wi:9 Ch:15        
        Neutral $:993 HP:9(12) Pw:3(3) AC:10 Exp:1/19 Hungry          
        Dlvl:1 T:752                                  Conf            
\end{verbatim}
\begin{center}
Figure 2
\end{center}
}

%.hn 2
\subsection*{The status lines (bottom)}

The bottom two (or three) lines of the screen contain several cryptic
pieces of information describing your current status.
Figure 1 shows the traditional two-line status area below the map.
Figure 2 shows just the status area, when the \textit{statuslines:3}
option has been set (not all interfaces support this option).
If any status line becomes wider than the screen, you might not see all
of it due to truncation.
When the numbers grow bigger and multiple \textit{conditions} are present,
the two-line format will run out of room on the second line, but
\textit{statuslines:2}
is the default because a basic 24-line terminal isn't tall enough for
the third line.

%.pg
Here are explanations of what the various status items mean:

%.lp
\begin{description}
\item[Title]
Your character's name and professional ranking (based on role and
\textit{experience level}, see below).
%.lp
\item[Strength]
A measure of your character's strength; one of your six basic
attributes.  A human character's attributes can range from 3 to 18 inclusive;
non-humans may exceed these limits
(occasionally you may get super-strengths of the form 18/xx, and magic can
also cause attributes to exceed the normal limits).  The
higher your strength, the stronger you are.  Strength affects how
successfully you perform physical tasks, how much damage you do in
combat, and how much loot you can carry.
%.lp
\item[Dexterity]
Dexterity affects your chances to hit in combat, to avoid traps, and
do other tasks requiring agility or manipulation of objects.
%.lp
\item[Constitution]
Constitution affects your ability to recover from injuries and other
strains on your stamina.
When strength is low or modest, constitution also affects how much you
can carry.  With sufficiently high strength, the contribution to
carrying capacity from your constitution no longer matters.
%.lp
\item[Intelligence]
Intelligence affects your ability to cast spells and read spellbooks.
%.lp
\item[Wisdom]
Wisdom comes from your practical experience (especially when dealing with
magic).  It affects your magical energy.
%.lp
\item[Charisma]
Charisma affects how certain creatures react toward you.  In
particular, it can affect the prices shopkeepers offer you.
%.lp
\item[Alignment]
%
\textit{Lawful}, \textit{Neutral} or \textit{Chaotic}.  Often, Lawful is
taken as good and Chaotic as evil, but legal and ethical do not always
coincide.  Your alignment influences how other
monsters react toward you.  Monsters of a like alignment are more likely
to be non-aggressive, while those of an opposing alignment are more likely
to be seriously offended at your presence.
%.lp
\item[Dungeon Level]
How deep you are in the dungeon.  You start at level one and the number
increases as you go deeper into the dungeon.  Some levels are special,
and are identified by a name and not a number.  The Amulet of Yendor is
reputed to be somewhere beneath the twentieth level.
%.lp
\item[Gold]
The number of gold pieces you are openly carrying.  Gold which you have
concealed in containers is not counted.
%.lp
\item[Hit Points]
Your current and maximum hit points.  Hit points indicate how much
damage you can take before you die.  The more you get hit in a fight,
the lower they get.  You can regain hit points by resting, or by using
certain magical items or spells.  The number in parentheses is the maximum
number your hit points can reach.
%.lp
\item[Power]
Spell points.  This tells you how much mystic energy (\textit{mana})
you have available for spell casting.  Again, resting will regenerate the
amount available.
%.lp
\item[Armor Class]
A measure of how effectively your armor stops blows from unfriendly
creatures.  The lower this number is, the more effective the armor; it
is quite possible to have negative armor class.
See the \textit{Armor} subsection of \textit{Objects} for more information.
%.lp
\item[Experience]
Your current experience level.
If the \textit{showexp}
option is set, it will be followed by a slash and experience points.
As you adventure, you gain experience points.
At certain experience point totals, you gain an experience level.
The more experienced you are, the better you fight and withstand magical
attacks.
(By the time your level reaches double digits, the usefulness of showing
the points with it has dropped significantly.
You can use the `\texttt{O}' command to turn \textit{showexp}
off to avoid using up the limited status line space.)
%.lp
\item[Time]
The number of turns elapsed so far, displayed if you have the
\textit{time} option set.
%.lp
\item[Status]
Hunger:
your current hunger status.
Values are \textit{Satiated}, \textit{Not~Hungry} (or \textit{Normal}),
\textit{Hungry}, \textit{Weak}, and \textit{Fainting}.
%.\" not mentioned: Fainted
Not shown when \textit{Normal}.

%.lp ""
Encumbrance:
an indication of how what you are carrying affects your ability to move.
Values are \textit{Unencumbered}, \textit{Burdened}, \textit{Stressed},
\textit{Strained}, \textit{Overtaxed}, and \textit{Overloaded}.
Not shown when \textit{Unencumbered}.

%.lp ""
Fatal~conditions:
\textit{Stone} (aka \textit{Petrifying}, turning to stone),
\textit{Slime} (turning into green slime),
\textit{Strngl} (being strangled),
\textit{FoodPois} (suffering from acute food poisoning),
\textit{TermIll} (suffering from a terminal illness).

%.lp ""
Non-fatal~conditions:
\textit{Blind} (can't see), \textit{Deaf} (can't hear),
\textit{Stun} (stunned), \textit{Conf} (confused), \textit{Hallu} (hallucinating).

%.lp ""
Movement~modifiers:
\textit{Lev} (levitating), \textit{Fly} (flying), \textit{Ride} (riding).

%.lp ""
Other conditions and modifiers exist, but there isn't enough room to
display them with the other status fields.
\\
% unindented paragraph
The \texttt{\#attributes} command (default key \texttt{\textasciicircum X}) will show
all current status information in unabbreviated format.
It also shows other information which might be included on the status
lines if those had more room.

\end{description}

%.hn 2
\subsection*{The message line (top)}

%.pg
The top line of the screen is reserved for messages that describe
things that are impossible to represent visually.  If you see a
``\texttt{--More--}'' on the top line, this means that \textit{NetHack} has
another message to display on the screen, but it wants to make certain
that you've read the one that is there first.  To read the next message,
just press the space bar.

%.pg
To change how and what messages are shown on the message line,
see ``\textit{Configuring Message Types}`` and the \textit{verbose}
option.

%.hn 2
\subsection*{The map (rest of the screen)}

%.pg
The rest of the screen is the map of the level as you have explored it
so far.  Each symbol on the screen represents something.  You can set
various graphics
options to change some of the symbols the game uses; otherwise, the
game will use default symbols.  Here is a list of what the default
symbols mean:

\begin{description}[font=\mdseries\ttfamily]
\item[-]
The horizontal or corner walls of a room, or an open east/west door.
\item[|]
The vertical walls of a room, or an open north/south door, or a grave.
\item[.]
The floor of a room, or ice, or a doorless doorway, or the span of an
open drawbridge.
\item[\#]
A corridor, or iron bars, or a tree, or the portcullis of a closed
drawbridge.\\
%.lp ""
Note: engravings in corridors also appear as \# but are shown in
a different color from normal corridor locations.
\item[>]
Stairs down: a way to the next level.
\item[<]
Stairs up: a way to the previous level.
\item[+]
A closed door, or a spellbook containing a spell you may be able to learn.
\item[@]
Your character or a human or an elf.
\item[\textdollar]
A pile of gold.
\item[\textasciicircum]
A trap (once you have detected it).
\item[)]
A weapon.
\item[{[}]
A suit or piece of armor.
\item[\%]
Something edible (not necessarily healthy).
\item[?]
A scroll.
\item[/]
A wand.
\item[=]
A ring.
\item[!]
A potion.
\item[(]
A useful item (pick-axe, key, lamp \ldots).
\item["]
An amulet or a spider web.
\item[*]
A gem or rock (possibly valuable, possibly worthless).
\item[\textasciigrave]
A boulder or statue or an engraving on the floor of a room.\\
%.lp ""
Note: statues are displayed as if they were the monsters they depict
so won't appear as a \textit{grave accent} (aka \textit{back-tick}).
\item[0]
An iron ball.
\item[\textunderscore]
An altar, or an iron chain.
\item[\textbraceleft]
A fountain or a sink.
\item[\textbraceright]
A pool of water or moat or a wall of water
or a pool of lava or a wall of lava.
\item[\textbackslash]
An opulent throne.
\item[a-z \textrm{\textmd{and}}]
\item[A-HJ-Z \textrm{\textmd{and}}]
\item[@\&\textquotesingle :;]
Letters and certain other symbols represent the various inhabitants
of the Mazes of Menace.
Watch out, they can be nasty and vicious.
Sometimes, however, they can be helpful.
\item[I]
Rather than a specific type of monster, this marks the last known
location of an invisible or otherwise unseen monster.
Note that the monster could have moved.
The `\texttt{s}', `\texttt{F}', and `\texttt{m}' commands may be useful here.
\item[1-5]
The digits 1 through 5 may be displayed, marking unseen monsters sensed
via the \textit{Warning} attribute.
Less dangerous monsters are indicated by lower values, more dangerous by
higher values.

\end{description}
%.pg
You need not memorize all these symbols; you can ask the game what any
symbol represents with the `\texttt{/}' command (see the next section for
more info).

%.hn 1
\section{Commands}

%.pg
Commands can be initiated by typing one or two characters to which
the command is bound to, or typing the command name in the extended
commands entry.  Some commands,
like ``\texttt{search}'', do not require that any more information be collected
by \textit{NetHack}.  Other commands might require additional information, for
example a direction, or an object to be used.  For those commands that
require additional information, \textit{NetHack} will present you with either
a menu of choices, or with a command line prompt requesting information.
Which you are presented with will depend chiefly on how you have set the
`\textit{menustyle}'
option.

%.pg
For example, a common question in the form ``\texttt{What do you want to
use? [a-zA-Z\ ?*]}'', asks you to choose an object you are carrying.
Here, ``\texttt{a-zA-Z}'' are the inventory letters of your possible choices.
Typing `\texttt{?}' gives you an inventory list of these items, so you can see
what each letter refers to.  In this example, there is also a `\texttt{*}'
indicating that you may choose an object not on the list, if you
wanted to use something unexpected.  Typing a `\texttt{*}' lists your entire
inventory, so you can see the inventory letters of every object you're
carrying.  Finally, if you change your mind and decide you don't want
to do this command after all, you can press the `ESC' key to abort the
command.

%.pg
You can put a number before some commands to repeat them that many
times; for example, ``\texttt{10s}'' will search ten times.  If you have the
\textit{number\textunderscore pad}
option set, you must type `\texttt{n}' to prefix a count, so the example above
would be typed ``\texttt{n10s}'' instead.  Commands for which counts make no
sense ignore them.  In addition, movement commands can be prefixed for
greater control (see below).  To cancel a count or a prefix, press the
`ESC' key.

%.pg
The list of commands is rather long, but it can be read at any time
during the game through the `\texttt{?}' command, which accesses a menu of
helpful texts.  Here are the default key bindings for your reference:

\begin{description}[font=\mdseries\ttfamily]
%.lp
\item[?]
Help menu:  display one of several help texts available.
%.lp
\item[/]
The \texttt{whatis} command, to
tell what a symbol represents.  You may choose to specify a location
or type a symbol (or even a whole word) to explain.
Specifying a location is done by moving the cursor to a particular spot
on the map and then pressing one of `\texttt{.}', `\texttt{,}', `\texttt{;}',
or `\texttt{:}'.  `\texttt{.}' will explain the symbol at the chosen location,
conditionally check for ``\texttt{More info?}'' depending upon whether the
`\textit{help}'
option is on, and then you will be asked to pick another location;
`\texttt{,}' will explain the symbol but skip any additional
information, then let you pick another location;
`\texttt{;}' will skip additional info and also not bother asking
you to choose another location to examine; `\texttt{:}' will show additional
info, if any, without asking for confirmation.  When picking a location,
pressing the \texttt{ESC} key will terminate this command, or pressing `\texttt{?}'
will give a brief reminder about how it works.

%.lp ""
If the
\textit{autodescribe}
option is on, a short description of what you see at each location is
shown as you move the cursor.  Typing `\texttt{\#}' while picking a location will
toggle that option on or off.
The
\textit{whatis\textunderscore coord}
option controls whether the short description includes map coordinates.

%.lp ""
Specifying a name rather than a location
always gives any additional information available about that name.

%.lp ""
You may also request a description of nearby monsters,
all monsters currently displayed, nearby objects, or all objects.
The
\textit{whatis\textunderscore coord}
option controls which format of map coordinate is included with their
descriptions.
%.lp
\item[\&]
Tell what a command does.
%.lp
\item[<]
Go up to the previous level (if you are on a staircase or ladder).
%.lp
\item[>]
Go down to the next level (if you are on a staircase or ladder).
%.lp
\item[{[yuhjklbn]}]
Go one step in the direction indicated (see Figure 3).  If you sense
or remember
a monster there, you will fight the monster instead.  Only these
one-step movement commands cause you to fight monsters; the others
(below) are ``safe.''
%.sd
\begin{center}
\begin{tabular}{cc}
\verb+   y  k  u   + & \verb+   7  8  9   +\\
\verb+    \ | /    + & \verb+    \ | /    +\\
\verb+   h- . -l   + & \verb+   4- . -6   +\\
\verb+    / | \    + & \verb+    / | \    +\\
\verb+   b  j  n   + & \verb+   1  2  3   +\\
                     & (if \textit{number\textunderscore pad} set)
\end{tabular}
\end{center}
%.ed
\begin{center}
Figure 3
\end{center}
%.lp
\item[{[YUHJKLBN]}]
Go in that direction until you hit a wall or run into something.
%.lp
\item[m{[yuhjklbn]}]
Prefix:  move without picking up objects or fighting (even if you remember
a monster there).\\
%.lp ""
A few non-movement commands use the `\texttt{m}' prefix to request
operating via menu (to temporarily override the
\textit{menustyle:Traditional}
option).
Primarily useful for `\texttt{,}' (pickup) when there is only one class of
objects present (where there won't be any ``what kinds of objects?'' prompt,
so no opportunity to answer `\texttt{m}' at that prompt).
\\
%.lp ""
The prefix will
make ``\texttt{\#travel}'' command show a menu of interesting targets in sight.
It can also be used with the `\texttt{\textbackslash}' (known, show a
list of all discovered objects) and the `\texttt{\`{}}' (knownclass,
show a list of discovered objects in a particular class) commands to offer
a menu of several sorting alternatives (which sets a new value for the
\textit{sortdiscoveries}
option); also for ``\texttt{\#vanquished}'' and ``\texttt{\#genocided}'' commands
to offer a sorting menu.
\\
%.lp ""
A few other commands (eat food, offer sacrifice, apply tinning-kit,
drink/quaff, dip, tip container) use
the `\texttt{m}' prefix to skip checking for applicable objects on
the floor and go straight to checking inventory,
or (for ``\texttt{\#loot}'' to remove a saddle),
skip containers and go straight to adjacent monsters.
\\
%.lp ""
In debug mode (aka ``wizard mode''), the `\texttt{m}' prefix may also be
used with the ``\texttt{\#teleport}'' and ``\texttt{\#wizlevelport}'' commands.
%.lp
\item[F{[yuhjklbn]}]
Prefix:  fight a monster (even if you only guess one is there).
%.lp
\item[g{[yuhjklbn]}]
Prefix:  Move until something interesting is found.
%.lp
\item[G{[yuhjklbn]} \textrm{\textmd{or}} <Control>+{[yuhjklbn]}]
Prefix:  Similar to `\texttt{g}', but forking of corridors is not considered
interesting.
\\
Note:  \texttt{<Control>+<key>} means holding the \texttt{<Control>} or
\texttt{<Ctrl>} key down like \texttt{<Shift>} while typing and releasing
\texttt{<key>}, then releasing \texttt{<Control>}.  \texttt{\textasciicircum <key>} is used as
shorthand elsewhere in the Guidebook to mean the same thing.  Control
characters are case-insensitive so \texttt{\textasciicircum x} and \texttt{\textasciicircum X} are the same.
%.lp
\item[M{[yuhjklbn]}]
Old versions supported `\texttt{M}' as a movement prefix which
combined the effect of `\texttt{m}' with \texttt{<Control>+<direction>}.
That is no longer supported as a prefix but similar effect can be achieved
by using \texttt{m} and \texttt{G<direction>} in combination.
\texttt{m} can also be used in combination with \texttt{g<direction>},
\texttt{<Control>+<direction>}, or \texttt{<Shift>+<direction>}.
%.lp
\item[\textunderscore]
Travel to a map location via a shortest-path algorithm.\\
%.lp ""
The shortest path
is computed over map locations the hero knows about (e.g. seen or
previously traversed).
If there is no known path, a guess is made instead.
Stops on most of
the same conditions as the `\texttt{G}' command, but without picking up
objects, so implicitly forces the `\texttt{m}' prefix.
For ports with mouse
support, the command is also invoked when a mouse-click takes place on a
location other than the current position.
%.lp
\item[.]
Wait or rest, do nothing for one turn.
Precede with the `\texttt{m}' prefix
to wait for a turn even next to a hostile monster, if \textit{safe\textunderscore wait}
is on.
%.lp
\item[a]
Apply (use) a tool (pick-axe, key, lamp \ldots).\\
%.lp ""
If used on a wand, that wand will be broken, releasing its magic in the
process.
Confirmation is required.
%.lp
\item[A]
Remove one or more worn items, such as armor.\\
%.lp ""
Use `\texttt{T}' (take off) to take off only one piece of armor
or `\texttt{R}' (remove) to take off only one accessory.
%.lp
\item[\textasciicircum A]
Repeat the previous command.
%.lp
\item[c]
Close a door.
%.lp
\item[C]
Call (name) a monster, an individual object, or a type of object.\\
%.lp ""
Same as extended command ``\texttt{\#name}''.
%.lp
\item[\textasciicircum C]
Panic button.  Quit the game.
%.lp
\item[d]
Drop something.\\
For example \texttt{d7a} --- drop seven items of object
\textit{a}.
%.lp
\item[D]
Drop several things.\\
%.lp ""
In answer to the question\\
``\texttt{What kinds of things do you want to drop? [!\%= BUCXPaium]}''\\
you should type zero or more object symbols possibly followed by
`\texttt{a}' and/or `\texttt{i}' and/or `\texttt{u}' and/or `\texttt{m}'.
In addition, one or more of
the bless\-ed/\-un\-curs\-ed/\-curs\-ed groups may be typed.\\
%.sd
%.si
\texttt{DB}  --- drop all objects known to be blessed.\\
\texttt{DU}  --- drop all objects known to be uncursed.\\
\texttt{DC}  --- drop all objects known to be cursed.\\
\texttt{DX}  --- drop all objects of unknown B/U/C status.\\
\texttt{DP}  --- drop objects picked up last.\\
\texttt{Da}  --- drop all objects, without asking for confirmation.\\
\texttt{Di}  --- examine your inventory before dropping anything.\\
\texttt{Du}  --- drop only unpaid objects (when in a shop).\\
\texttt{Dm}  --- use a menu to pick which object(s) to drop.\\
\texttt{D\%u} --- drop only unpaid food.
%.ei
%.ed
The last example shows a combination.
There are four categories of object filtering: class (`\texttt{!}' for
potions, `\texttt{?}' for scrolls, and so on), shop status (`\texttt{u}' for
unpaid, in other words, owned by the shop), bless/curse state
(`\texttt{B}', `\texttt{U}', `\texttt{C}', and `\texttt{X}' as shown above),
and novelty (`\texttt{P}', recently picked up items; controlled by picking
up or dropping things rather than by any time factor).
%.lp ""
\\
If you specify more than one value in a category (such as ``\texttt{!?}'' for
potions and scrolls or ``\texttt{BU}'' for blessed and uncursed), an inventory
object will meet the criteria if it matches any of the specified
values (so ``\texttt{!?}'' means `\texttt{!}' or `\texttt{?}').
If you specify more than one category, an inventory object must meet
each of the category criteria (so ``\texttt{\%u}'' means class `\texttt{\%}' and
unpaid `\texttt{u}').
Lastly, you may specify multiple values within multiple categories:
``\texttt{!?BU}'' will select all potions and scrolls which are known to be
blessed or uncursed.
(In versions prior to 3.6, filter combinations behaved differently.)
%.lp
\item[\textasciicircum D]
Kick something (usually a door).
%.lp
\item[e]
Eat food.\\
%.lp ""
Normally checks for edible item(s) on the floor, then if none are found
or none are chosen, checks for edible item(s) in inventory.
Precede `\texttt{e}' with the `\texttt{m}' prefix to bypass attempting to eat
anything off the floor.\\
%.lp ""
If you attempt to eat while already satiated, you might choke to death.
If you risk it, you will be asked whether
to ``continue eating?'' \textit{if you survive the first bite}.
You can set the
\textit{paranoid\textunderscore confirmation:eating}
option to require a response of ``\texttt{yes}'' instead of just `\texttt{y}'.
%.lp
% Make sure Elbereth is not hyphenated below, the exact spelling matters.
% (Only specified here to parallel Guidebook.mn; use of \tt font implicitly
% prevents automatic hyphenation in TeX and LaTeX.)
\hyphenation{Elbereth}		%override the deduced syllable breaks
\item[E]
Engrave a message on the floor.\\
%.sd
%.si
\texttt{E-} --- write in the dust with your fingers.\\
%.ei
%.ed
%.lp ""
Engraving the word ``\texttt{Elbereth}'' will cause most monsters to not attack
you hand-to-hand (but if you attack, you will rub it out); this is
often useful to give yourself a breather.
%.lp
\item[f]
Fire (shoot or throw) one of the objects placed in your quiver (or
quiver sack, or that you have at the ready).
You may select ammunition with a previous `\texttt{Q}' command, or let the
computer pick something appropriate if \textit{autoquiver} is true.
If your wielded weapon has the throw-and-return property, your quiver
is empty, and \textit{autoquiver}
is false, you will throw that wielded weapon instead of filling the quiver.
This will also automatically use a polearm if wielded.
If \textit{fireassist} is true, firing will automatically try to wield a launcher
(for example, a bow or a sling) matching the ammo in the quiver; this might
take multiple turns, and get interrupted by a monster.
Remember to swap back to your main melee weapon afterwards.
%.lp ""
\\
See also `\texttt{t}' (throw) for more general throwing and shooting.
%.lp
\item[i]
List your inventory (everything you're carrying).
%.lp
\item[I]
List selected parts of your inventory, usually be specifying the character
for a particular set of objects, like `\texttt{[}' for armor or `\texttt{!}'
for potions.\\
%.sd
%.si
\texttt{I*} --- list all gems in inventory;\\
\texttt{Iu} --- list all unpaid items;\\
\texttt{Ix} --- list all used up items that are on your shopping bill;\\
\texttt{IB} --- list all items known to be blessed;\\
\texttt{IU} --- list all items known to be uncursed;\\
\texttt{IC} --- list all items known to be cursed;\\
\texttt{IX} --- list all items whose bless/curse status is unknown;\\
\texttt{IP} --- list items picked up last;\\
\texttt{I\$} --- count your money.
%.ei
%.ed
%.lp
\item[o]
Open a door.
%.lp
\item[O]
Set options.\\
%.lp ""
A menu showing the current option values will be
displayed.  You can change most values simply by selecting the menu
entry for the given option (ie, by typing its letter or clicking upon
it, depending on your user interface).  For the non-boolean choices,
a further menu or prompt will appear once you've closed this menu.
The available options
are listed later in this Guidebook.  Options are usually set before the
game rather than with the `\texttt{O}' command; see the section on options below.
Precede \texttt{O} with the \texttt{m} prefix to show advanced options.
%.lp
\item[\textasciicircum O]
Show overview.\\
%.lp ""
Shortcut for ``\texttt{\#overview}'':
list interesting dungeon levels visited.\\
%.lp ""
(Prior to 3.6.0, `\texttt{\textasciicircum O}' was a debug mode command which listed
the placement of all special levels.
Use ``\texttt{\#wizwhere}'' to run that command.)
%.lp
\item[p]
Pay your shopping bill.
%.lp
\item[P]
Put on an accessory (ring, amulet, or blindfold).\\
%.lp ""
This command may also be used to wear armor.  The prompt for
which inventory item to use will only list accessories, but choosing
an unlisted item of armor will attempt to wear it.
(See the `\texttt{W}' command below.  It lists armor as the inventory
choices but will accept an accessory and attempt to put that on.)
%.lp
\item[\textasciicircum P]
Repeat previous message.\\
%.lp ""
Subsequent \texttt{\textasciicircum P}'s repeat earlier messages.
For some interfaces, the behavior can be varied via the
\textit{msg\textunderscore window} option.
%.lp
\item[q]
Quaff (drink) something (potion, water, etc).\\
%.lp ""
When there is a fountain or sink present, it asks whether to drink
from that.
If that is declined, then it offers a chance to choose a potion from
inventory.
Precede \texttt{q} with the \texttt{m} prefix to skip asking about
drinking from a fountain or sink.
%.lp
\item[Q]
Select an object for your quiver, quiver sack, or just generally at
the ready (only one of these is available at a time).  You can then throw
this (or one of these) using the `\texttt{f}' command.
%.lp
\item[r]
Read a scroll or spellbook.
%.lp
\item[R]
Remove a worn accessory (ring, amulet, or blindfold).\\
%.lp ""
If you're wearing more than one, you'll be prompted for which one to
remove.  When you're only wearing one, then by default it will be removed
without asking, but you can set the
\textit{paranoid\textunderscore confirmation:Remove}
option to require a prompt.\\
%.lp ""
This command may also be used to take off armor.  The prompt for which
inventory item to remove only lists worn accessories, but an item of
worn armor can be chosen.
(See the `\texttt{T}' command below.  It lists armor as the inventory
choices but will accept an accessory and attempt to remove it.)
%.lp
\item[\textasciicircum R]
Redraw the screen.
%.lp
\item[s]
Search for secret doors and traps around you.
It usually takes several tries to find something.
Precede with the `\texttt{m}' prefix to wait for a turn
even next to a hostile monster, if \textit{safe\textunderscore wait}
is on.\\
%.lp ""
Can also be used to figure out whether there is still a monster at
an adjacent ``remembered, unseen monster'' marker.
%.lp
\item[S]
Save the game (which suspends play and exits the program).
The saved game will be restored automatically the next time you play
using the same character name.\\
%.lp ""
In normal play, once a saved game is restored the file used to hold
the saved data is deleted.
In explore mode, once restoration is accomplished you are asked whether
to keep or delete the file.
Keeping the file makes it feasible to play for a while then quit
without saving and later restore again.\\
%.lp ""
There is no ``save current game state and keep playing'' command, not
even in explore mode where saved game files can be kept and re-used.
%.lp
\item[t]
Throw an object or shoot a projectile.\\
%.lp ""
There's no separate ``shoot'' command.
If you ``throw'' an arrow while wielding a bow, you are shooting
that arrow and any weapon skill bonus or penalty for bow applies.
If you ``throw'' an arrow while not wielding a bow, you are throwing
it by hand and it will generally be less effective than when shot.\\
%.lp ""
See also `\texttt{f}' (fire) for throwing or shooting an item pre-selected
via the `\texttt{Q}' (quiver) command, with some extra assistance.
%.lp
\item[T]
Take off armor.\\
%.lp ""
If you're wearing more than one piece, you'll be prompted for which
one to take off.  (Note that this treats a cloak covering a suit
and/or a shirt, or a suit covering a shirt, as if the underlying items
weren't there.)
When you're only wearing one, then by default it will
be taken off without asking, but you can set the
\textit{paranoid\textunderscore confirmation:Remove}
option to require a prompt.\\
%.lp ""
This command may also be used to remove accessories.  The prompt
for which inventory item to take off only lists worn armor, but a worn
accessory can be chosen.
(See the `\texttt{R}' command above.  It lists accessories as the inventory
choices but will accept an item of armor and attempt to take it off.)
%.lp
\item[\textasciicircum T]
Teleport, if you have the ability.
%.lp
\item[v]
Display a list of significant events in the current game, also shown by
\texttt{\#chronicle}.
\\
%.lp ""
`\texttt{v}' used to display the game's version number.
Use `\texttt{V}' for that now.

%.lp
\item[V]
Display version number.
\\
%.lp ""
`\texttt{V}' used to display a summary of the game's history.
Still available via \texttt{\#history}.
%.lp
\item[w]
Wield weapon.\\
%.sd
%.si
\texttt{w-} --- wield nothing, use your bare (or gloved) hands.\\
%.ei
%.ed
Some characters can wield two weapons at once; use the `\texttt{X}' command
(or the ``\texttt{\#twoweapon}'' extended command) to do so.
%.lp
\item[W]
Wear armor.\\
%.lp ""
This command may also be used to put on an accessory (ring, amulet, or
blindfold).  The prompt for which inventory item to use will only list
armor, but choosing an unlisted accessory will attempt to put it on.
(See the `\texttt{P}' command above.  It lists accessories as the inventory
choices but will accept an item of armor and attempt to wear it.)
%.lp
\item[x]
Exchange your wielded weapon with the item in your alternate weapon slot.\\
%.lp ""
The latter is used as your secondary weapon when engaging in
two-weapon combat.  Note that if one of these slots is empty,
the exchange still takes place.
%.lp
\item[X]
Toggle two-weapon combat, if your character can do it.  Also available
via the ``\texttt{\#twoweapon}'' extended command.\\
%.lp ""
(In versions prior to 3.6 this keystroke ran the command to switch from normal
play to ``explore mode'', also known as ``discovery mode'', which has now
been moved to ``\texttt{\#exploremode}'' and \texttt{M-X}.)
%.lp
\item[\textasciicircum X]
Display basic information about your character.\\
%.lp ""
Displays name, role, race, gender (unless role name makes that
redundant, such as \texttt{Caveman} or \texttt{Priestess}), and alignment,
along with your patron deity and his or her opposition.  It also
shows most of the various items of information from the status line(s)
in a less terse form, including several additional things which don't
appear in the normal status display due to space considerations.\\
%.lp ""
In normal play, that's all that `\texttt{\textasciicircum X}' displays.
In explore mode, the role and status feedback is augmented by the
information provided by \textit{enlightenment} magic.
%.lp
\item[z]
Zap a wand.\\
%.sd
%.si
\texttt{z.} --- to aim at yourself, use `\texttt{.}' for the direction.
%.ei
%.ed
%.lp
\item[Z]
Zap (cast) a spell.\\
%.sd
%.si
\texttt{Z.} --- to cast at yourself, use `\texttt{.}' for the direction.
%.ei
%.ed
%.lp
\item[\textasciicircum Z]
Suspend the game (UNIX versions with job control only).
See ``\#suspend'' below for more details.
%.lp
\item[:]
Look at what is here.
%.lp
\item[;]
Show what type of thing a visible symbol corresponds to.
%.lp
\item[,]
Pick up some things from the floor beneath you.\\
%.lp ""
May be preceded by `\texttt{m}' to force a selection menu.
%.lp
\item[@]
Toggle the \textit{autopickup} option on and off.
%.lp
\item[\textasciicircum]
Ask for the type of an adjacent trap you found earlier.
%.lp
\item[)]
Tell what weapon you are wielding.
%.lp
\item[{]}]
Tell what armor you are wearing.
%.lp
\item[=]
Tell what rings you are wearing.
%.lp
\item["]
Tell what amulet you are wearing.
%.lp
\item[(]
Tell what tools you are using.
%.lp
\item[*]
Tell what equipment you are using.\\
%.lp ""
Combines the preceding five type-specific
commands into one.
%.lp
\item[\$]
Report the gold you're carrying, possibly shop credit and/or debt too.
%.lp
\item[+]
List the spells you know.\\
%.lp ""
Using this command, you can also rearrange
the order in which your spells are listed, either by sorting the entire
list or by picking one spell from the menu then picking another to swap
places with it.  Swapping pairs of spells changes their casting letters,
so the change lasts after the current `\texttt{+}' command finishes.  Sorting
the whole list is temporary.  To make the most recent sort order persist
beyond the current `\texttt{+}' command, choose the sort option again and then
pick ``reassign casting letters''.  (Any spells learned after that will
be added to the end of the list rather than be inserted into the sorted
ordering.)
%.lp
\item[\textbackslash]
Show what types of objects have been discovered.
\\
%.lp ""
May be preceded by `\texttt{m}' to select preferred display order.
%.lp
\item[\textasciigrave]
Show discovered types for one class of objects.
\\
%.lp ""
May be preceded by `\texttt{m}' to select preferred display order.

%.lp
\item[|]
If persistent inventory display is supported and enabled (with the
\textit{perm\textunderscore invent}
option), interact with it instead of with the map.
\\
%.lp ""
Allows scrolling with the
menu\textunderscore first\textunderscore page, menu\textunderscore previous\textunderscore page,
menu\textunderscore next\textunderscore page, and menu\textunderscore last\textunderscore page
keys (`\texttt{\textasciicircum}', `\texttt{<}', `\texttt{>}', `\texttt{|}' by default).
Some interfaces also support menu\textunderscore shift\textunderscore left and menu\textunderscore shift\textunderscore right
keys (`\texttt{\textbraceleft}' and `\texttt{\textbraceright}' by default).
Use the \textit{Return} (aka \textit{Enter}) or \textit{Escape} key to
resume play.

%.lp
\item[!]
Escape to a shell.
See ``\#shell'' below for more details.
%.lp
\item[Del]
Show map without obstructions.
You can view the explored portion of the current level's map without
monsters; without monsters and objects; or without monsters, objects,
and traps.\\
%.lp ""
The \texttt{<del>} key is also shown as \texttt{<delete>} on some keyboards or
\texttt{<rubout>} on others.
It is sometimes displayed as \texttt{\textasciicircum ?} even though that is not an actual
control character.\\
%.lp ""
Many terminals have an option to swap the \texttt{<delete>} and \texttt{<backspace>}
keys, so typing the \texttt{<del>} key might not execute this command.
If that happens, you can use the extended command ``\texttt{\#terrain}'' instead.
%.lp
\item[\#]
Perform an extended command.\\
%.lp ""
As you can see, the authors of \textit{NetHack}
used up all the letters, so this is a way to introduce the less frequently
used commands.
What extended commands are available depends on what features
the game was compiled with.
%.lp
\item[\#adjust]
Adjust inventory letters (most useful when the
\textit{fixinv}
option is ``on''). Autocompletes. Default key is `\texttt{M-a}'.\\
%.lp ""
This command allows you to move an item from one particular inventory
slot to another so that it has a letter which is more meaningful for you
or that it will appear in a particular location when inventory listings
are displayed.
You can move to a currently empty slot, or if the destination is
occupied---and won't merge---the
item there will swap slots with the one being moved.
``\texttt{\#adjust}'' can also be used to split a stack of objects; when
choosing the item to adjust, enter a count prior to its letter.\\
%.lp ""
Adjusting without a count used to collect all compatible stacks when
moving to the destination.  That behavior has been changed; to gather
compatible stacks, ``\texttt{\#adjust}'' a stack into its own inventory slot.
If it has a name assigned, other stacks with the same name or with
no name will merge provided that all their other attributes match.
If it does not have a name, only other stacks with no name are eligible.
In either case, otherwise compatible stacks with a different name
will not be merged.  This contrasts with using ``\texttt{\#adjust}'' to move
from one slot to a different slot.  In that situation, moving (no
count given) a compatible stack will merge if either stack has a
name when the other doesn't and give that name to the result, while
splitting (count given) will ignore the source stack's name when
deciding whether to merge with the destination stack.
%.lp
\item[\#annotate]
Allows you to specify one line of text to associate with the current
dungeon level.  All levels with annotations are displayed by the
``\texttt{\#overview}'' command. Autocompletes.
Default key is `\texttt{M-A}',
and also `\texttt{\textasciicircum N}' if \textit{number\textunderscore pad} is on.
\\
%.lp ""
Preceding \#annotate with the `\texttt{m}' prefix is the same as
\#overview with the prefix.
%.lp
\item[\#apply]
Apply (use) a tool such as a pick-axe, a key, or a lamp.
Default key is `\texttt{a}'.\\
%.lp ""
If the tool used acts on items on the floor, using the `\texttt{m}' prefix
skips those items.\\
%.lp ""
If used on a wand, that wand will be broken, releasing its magic in the
process.  Confirmation is required.
%.lp
\item[\#attributes]
Show your attributes. Default key is `\texttt{\textasciicircum X}'.
%.lp
\item[\#autopickup]
Toggle the \textit{autopickup} option. Default key is `\texttt{@}'.
%.lp
\item[\#bugreport]
Bring up a browser window to submit a report to the \textit{NetHack Development
Team}.
Can be disabled at the time the program is built; when enabled,
CRASHREPORTURL must be set in the system configuration file.
%.lp
\item[\#call]
Call (name) a monster, or an object in inventory, on the floor,
or in the discoveries list, or add an annotation for the
current level (same as ``\texttt{\#annotate}''). Default key is `\texttt{C}'.
%.lp
\item[\#cast]
Cast a spell. Default key is `\texttt{Z}'.
%.lp
\item[\#chat]
Talk to someone. Default key is `\texttt{M-c}'.
%.lp
\item[\#chronicle]
Show a list of important game events. Default key is `\texttt{v}'.
%.lp
\item[\#close]
Close a door. Default key is `\texttt{c}'.
%.lp
\item[\#conduct]
List voluntary challenges you have maintained. Autocompletes.
Default key is `\texttt{M-C}'.\\
%.lp ""
See the section below entitled ``Conduct'' for details.
%.lp
\item[\#debugfuzzer]
Start the fuzz tester.
Debug mode only.
%.lp
\item[\#dip]
Dip an object into something. Autocompletes. Default key is `\texttt{M-d}'.\\
%.lp ""
The \texttt{m} prefix skips dipping into a fountain or pool if there
is one at your location.
%.lp
\item[\#down]
Go down a staircase. Default key is `\texttt{>}'.
%.lp
\item[\#drop]
Drop an item. Default key is `\texttt{d}'.
%.lp
\item[\#droptype]
Drop specific item types. Default key is `\texttt{D}'.
%.lp
\item[\#eat]
Eat something. Default key is `\texttt{e}'.
The `\texttt{m}' prefix skips eating items on the floor.
%.lp
\item[\#engrave]
Engrave writing on the floor. Default key is `\texttt{E}'.
%.lp
\item[\#enhance]
Advance or check weapon and spell skills. Autocompletes.
Default key is `\texttt{M-e}'.
%.lp
\item[\#exploremode]
Switch from normal play to non-scoring explore mode.
Default key is `\texttt{M-X}'.\\
%.lp ""
Requires confirmation; default response is `\texttt{n}' (no).
To really switch to explore mode, respond with `\texttt{y}'.
You can set the
\textit{paranoid\textunderscore confirmation:quit}
option to require a response of ``\texttt{yes}'' instead.
%.lp
\item[\#fight]
Prefix key to force fight a direction, even if you see nothing
to fight there.
Default key is `\texttt{F}', or `\texttt{-}' with
\textit{number\textunderscore pad}
%.lp
\item[\#fire]
Fire ammunition from quiver, possibly autowielding a launcher,
or hit with a wielded polearm.
Default key is `\texttt{f}'.
%.lp
\item[\#force]
Force a lock. Autocompletes. Default key is `\texttt{M-f}'.
%.lp
\item[\#genocided]
List any monster types which have been genocided.
In explore mode and debug mode it also shows types which have become
extinct.
\\
%.lp ""
The display order is the same as is used by \texttt{\#vanquished}.
The `\texttt{m}' prefix brings up a menu of available sorting orders, and
doing that for either \texttt{\#genocided} or \texttt{\#vanquished} changes the order for both.
\\
%.lp ""
If the sorting order is ``count high to low'' or ``count low to high''
(which are applicable for \texttt{\#vanquished}), that will be ignored
for \texttt{\#genocided} and alphabetical will be used instead.
The menu omits those two choices when used for \texttt{\#genocide}.
\\
%.lp ""
Autocompletes.
Default key is `\texttt{M-g}'.
%.lp
\item[\#glance]
Show what type of thing a map symbol corresponds to. Default key is `\texttt{;}'.
%.lp
\item[\#help]
Show the help menu.
Default key is `\texttt{?}',
and also `\texttt{h}' if \textit{number\textunderscore pad} is on.
%.lp
\item[\#herecmdmenu]
Show a menu of possible actions directed at your current location.
The menu is limited to a subset of the likeliest actions, not an
exhaustive set of all possibilities.
Autocompletes.\\
%.lp ""
If mouse support is enabled and the \textit{herecmd\textunderscore menu}
option is On, clicking on the hero (or steed when mounted) will
execute this command.
%.lp
\item[\#history]
Show a summary of the game's development.
%.lp
\item[\#inventory]
Show your inventory. Default key is `\texttt{i}'.
%.lp
\item[\#inventtype]
Inventory specific item types. Default key is `\texttt{I}'.
%.lp
\item[\#invoke]
Invoke an object's special powers. Autocompletes. Default key is `\texttt{M-i}'.
%.lp
\item[\#jump]
Jump to another location. Autocompletes.
Default key is `\texttt{M-j}',
and also `\texttt{j}' if \textit{number\textunderscore pad} is on.
%.lp
\item[\#kick]
Kick something.
Default key is `\texttt{\textasciicircum D}',
and also `\texttt{k}' if \textit{number\textunderscore pad} is on.
%.lp
\item[\#known]
Show what object types have been discovered.
Default key is `\texttt{\textbackslash}'.
\\
%.lp ""
The `\texttt{m}' prefix allows assigning a new value to the
\textit{sortdiscoveries}
option to control the order in which the discoveries are displayed.
%.lp
\item[\#knownclass]
Show discovered types for one class of objects.
Default key is `\texttt{`}'.
\\
%.lp ""
The `\texttt{m}' prefix operates the same as for \texttt{\#known}.
%.lp
\item[\#levelchange]
Change your experience level.
Autocompletes.
Debug mode only.
%.lp
\item[\#lightsources]
Show mobile light sources.
Autocompletes.
Debug mode only.
%.lp
\item[\#look]
Look at what is here, under you. Default key is `\texttt{:}'.
%.lp
\item[\#lookaround]
Describe what you can see, or remember, of your surroundings.
%.lp
\item[\#loot]
Loot a box or bag on the floor beneath you, or the saddle
from a steed standing next to you. Autocompletes.
Precede with the `\texttt{m}' prefix to skip containers at your location
and go directly to removing a saddle.
Default key is `\texttt{M-l}',
and also `\texttt{l}' if \textit{number\textunderscore pad} is on.
%.lp
\item[\#monster]
Use a monster's special ability (when polymorphed into monster form).
Autocompletes. Default key is `\texttt{M-m}'.
%.lp
\item[\#name]
Name a monster, an individual object, or a type of object.
Same as ``\texttt{\#call}''.
Autocompletes.
Default keys are `\texttt{N}', `\texttt{M-n}', and `\texttt{M-N}'.
%.lp
\item[\#offer]
Offer a sacrifice to the gods. Autocompletes. Default key is `\texttt{M-o}'.\\
%.lp ""
You'll need to find an altar to have any chance at success.
Corpses of recently killed monsters are the fodder of choice.
%.lp ""
The `\texttt{m}' prefix skips offering any items which are on the altar.\\
%.lp
\item[\#open]
Open a door. Default key is `\texttt{o}'.
%.lp
\item[\#options]
Show and change option settings. Default key is `\texttt{O}'.
Precede with the \texttt{m} prefix to show advanced options.
%.lp
\item[\#optionsfull]
Show advanced game option settings.
No default key.
Precede with the `\texttt{m}' prefix to execute the simpler options command.
(Mainly useful if you use \texttt{BINDING=O:optionsfull} to switch
`\texttt{O}' from simple options back to traditional advanced options.)
%.lp
\item[\#overview]
Display information you've discovered about the dungeon.
Any visited level
% [note: amnesia no longer causes levels to be forgotten so exclude this]
% (unless forgotten due to amnesia)
with an annotation is included,
and many things (altars, thrones, fountains, and so on; extra stairs
leading to another dungeon branch) trigger an automatic annotation.
If dungeon overview is chosen during end-of-game disclosure, every visited
level will be included regardless of annotations.
\\
%.lp ""
Precede \#overview with the `\texttt{m}' prefix to display the dungeon
overview as a menu where you can select any visited level to add or
remove an annotation without needing to return to that level.
This will also force all visited levels to be displayed rather than just
the ``interesting'' subset.
\\
%.lp ""
Autocompletes.
Default keys are `\texttt{\textasciicircum O}', and `\texttt{M-O}'.
% DON'T PANIC!
%.lp
\item[\#panic]
Test the panic routine.
Terminates the current game.
Autocompletes.
Debug mode only.\\
%.lp ""
Asks for confirmation; default is `\texttt{n}' (no); continue playing.
To really panic, respond with `\texttt{y}'.
You can set the
\textit{paranoid\textunderscore confirmation:quit}
option to require a response of ``\texttt{yes}'' instead.
%.lp
\item[\#pay]
Pay your shopping bill. Default key is `\texttt{p}'.
%.lp
\item[\#perminv]
If persistent inventory display is supported and enabled (with the
\textit{perm\textunderscore invent} option), interact with it instead of with the map.
You'll be prompted for menu scrolling keystrokes such
as `\texttt{>}' and `\texttt{<}'.
Press \texttt{Return} or \texttt{Escape} to resume normal play.
Default key is \texttt{|}.
%.lp
\item[\#pickup]
Pick up things at the current location. Default key is `\texttt{,}'.
The `\texttt{m}' prefix forces use of a menu.
%.lp
\item[\#polyself]
Polymorph self.
Autocompletes.
Debug mode only.
%.lp
\item[\#pray]
Pray to the gods for help. Autocompletes. Default key is `\texttt{M-p}'.\\
%.lp ""
Praying too soon after receiving prior help is a bad idea.
(Hint: entering the dungeon alive is treated as having received help.
You probably shouldn't start off a new game by praying right away.)
Since using this command by accident can cause trouble, there is an
option to make you confirm your intent before praying.  It is enabled
by default, and you can reset the
\textit{paranoid\textunderscore confirmation}
option to disable it.
%.lp
\item[\#prevmsg]
Show previously displayed game messages. Default key is `\texttt{\textasciicircum P}'.
%.lp
\item[\#puton]
Put on an accessory (ring, amulet, etc). Default key is `\texttt{P}'.
%.lp
\item[\#quaff]
Quaff (drink) something. Default key is `\texttt{q}'.\\
%.lp ""
The \texttt{m} prefix skips drinking from a fountain or sink if there
is one at your location.
%.lp
\item[\#quit]
Quit the program without saving your game. Autocompletes.\\
%.lp ""
Since using this command by accident would throw away the current game,
you are asked to confirm your intent before quitting.
Default response is `\texttt{n}' (no); continue playing.
To really quit, respond with `\texttt{y}'.
You can set the
\textit{paranoid\textunderscore confirmation:quit}
option to require a response of ``\texttt{yes}'' instead.
%.lp
\item[\#quiver]
Select ammunition for quiver. Default key is `\texttt{Q}'.
%.lp
\item[\#read]
Read a scroll, a spellbook, or something else. Default key is `\texttt{r}'.
%.lp
\item[\#redraw]
Redraw the screen.
Default key is `\texttt{\textasciicircum R}',
and also `\texttt{\textasciicircum L}' if \textit{number\textunderscore pad} is on.
%.lp
\item[\#remove]
Remove an accessory (ring, amulet, etc). Default key is `\texttt{R}'.
%.lp
\item[\#repeat]
Repeat the previous command.
Default key is~`\texttt{\textasciicircum A}'.
%.lp
\item[\#reqmenu]
Prefix key to modify the behavior or request menu from some commands.
Prevents autopickup when used with movement commands.
Default key is `\texttt{m}'.
%.lp
\item[\#retravel]
Travel to a previously selected travel destination.
Default key is `\texttt{C-\textunderscore }'.
See also \texttt{\#travel}.
%.lp
\item[\#ride]
Ride (or stop riding) a saddled creature. Autocompletes.
Default key is `\texttt{M-R}'.
%.lp
\item[\#rub]
Rub a lamp or a stone. Autocompletes. Default key is `\texttt{M-r}'.
%.lp
\item[\#run]
Prefix key to run towards a direction.
Default key is `\texttt{G}' when
\textit{number\textunderscore pad}
is off,
`\texttt{5}' when
\textit{number\textunderscore pad}
is set to 1~or~3,
otherwise `\texttt{M-5}' when it is set to 2~or~4.
%.lp
\item[\#rush]
Prefix key to rush towards a direction.
Default key is `\texttt{g}' when
\textit{number\textunderscore pad}
is off,
`\texttt{M-5}' when
\textit{number\textunderscore pad}
is set to 1~or~3,
otherwise `\texttt{5}' when it is set to 2~or~4.
%.lp
\item[\#save]
Save the game and exit the program.
Default key is `\texttt{S}'.
%.lp
\item[\#saveoptions]
Save configuration options to the config file.
This will overwrite the file, removing all comments, so if you have
manually edited the config file, don't use this.
%.lp
\item[\#search]
Search for traps and secret doors around you. Default key is `\texttt{s}'.
%.lp
\item[\#seeall]
Show all equipment in use. Default key is `\texttt{*}'.
%.lp ""
\\
Will display in-use items in a menu even when there is only one.
%.lp
\item[\#seeamulet]
Show the amulet currently worn. Default key is `\texttt{"}'.
%.lp ""
\\
Using the `\texttt{m}' prefix will force the display of a worn
amulet in a menu rather than with just a message.
%.lp
\item[\#seearmor]
Show the armor currently worn. Default key is `\texttt{[}'.
%.lp ""
\\
Will display worn armor in a menu even when there is only thing worn.
%.lp
\item[\#seerings]
Show the ring(s) currently worn. Default key is `\texttt{=}'.
%.lp ""
Will display worn rings in a menu if there are two (or there is
just one and is a meat ring rather than a ``real'' ring).
Use the `\texttt{m}' prefix to force a menu for one ring.
%.lp
\item[\#seetools]
Show the tools currently in use. Default key is `\texttt{(}'.
%.lp ""
Will display the result in a message if there is one tool in use (worn
blindfold or towel or lenses, lit lamp(s) and/or candle(s), leashes
attached to pets).
Will display a menu if there are more than one or if the command is
preceded by the `\texttt{m}' prefix.
%.lp
\item[\#seeweapon]
Show the weapon currently wielded. Default key is `\texttt{)}'.
%.lp ""
If dual-wielding, a separate message about the secondary weapon will be
given.
Using the `\texttt{m}' prefix will force a menu and it will include
primary weapon, alternate weapon even when not dual-wielding, and also
whatever is currently assigned to the quiver slot.
%.lp
\item[\#shell]
Do a shell escape, switching from NetHack to a subprocess.
Can be disabled at the time the program is built.
When enabled, access for specific users can be controlled by the system
configuration file.
Use the shell command `\texttt{exit}' to return to the game.
Default key is `\texttt{!}'.
%.lp
\item[\#showgold]
Report the gold in your inventory, including gold you know about in
containers you're carrying.  If you are inside a shop, report any credit
or debt you have in that shop.
Default key is `\texttt{\$}'.
%.lp
\item[\#showspells]
List and reorder known spells.
Default key is `\texttt{+}'.
%.lp
\item[\#showtrap]
Describe an adjacent trap, possibly covered by objects or a monster.
To be eligible, the trap must already be discovered.
(The ``\texttt{\#terrain}'' command can display your map with all objects and
monsters temporarily removed, making it possible to see all discovered
traps.)
Default key is `\texttt{\textasciicircum }'.
%.lp
\item[\#sit]
Sit down. Autocompletes. Default key is `\texttt{M-s}'.
%.lp
\item[\#stats]
Show memory usage statistics.
Autocompletes.
Debug mode only.
%.lp
\item[\#suspend]
Suspend the game, switching from NetHack to the terminal it was started
from without performing save-and-exit.
Can be disabled at the time the program is built.
When enabled, mainly useful for \textit{tty} and \textit{curses} interfaces on
%.UX \. \" yields "UNIX."
UNIX.
Use the shell command `\texttt{fg}' to return to the game.
Default key is `\texttt{\textasciicircum Z}'.
%.lp
\item[\#swap]
Swap wielded and secondary weapons. Default key is `\texttt{x}'.
%.lp
\item[\#takeoff]
Take off one piece of armor. Default key is `\texttt{T}'.
%.lp
\item[\#takeoffall]
Remove all armor. Default key is `\texttt{A}'.
%.lp
\item[\#teleport]
Teleport around the level. Default key is `\texttt{\textasciicircum T}'.
%.lp
\item[\#terrain]
Show map without obstructions.
In normal play you can view the explored portion of the current level's
map without monsters; without monsters and objects; or without monsters,
objects, and traps.\\
%.lp ""
If there are visible clouds of gas in view, they are treated like traps
when deciding whether to show them or the floor underneath them.\\
%.lp ""
In explore mode, you can choose to view the full map rather than just
its explored portion.
In debug mode there are additional choices.\\
%.lp ""
Autocompletes.
Default key is `\texttt{<del>}' or `\texttt{<delete>}' (see \textit{Del} above).
%.lp
\item[\#therecmdmenu]
Show a menu of possible actions directed at a location next to you.
The menu is limited to a subset of the likeliest actions, not an
exhaustive set of all possibilities.
Autocompletes.
%%--invoking it by mouse seems to be broken
%% \\
%% .lp ""
%% If mouse support is enabled and the \textit{herecmd\textunderscore menu}
%% option is On, clicking on an adjacent location will execute this command.
%.lp
\item[\#throw]
Throw something. Default key is `\texttt{t}'.
%.lp
\item[\#timeout]
Look at the timeout queue.
Autocompletes.
Debug mode only.
%.lp
\item[\#tip]
Tip over a container (bag or box) to pour out its contents.
When there are containers on the floor, the game will prompt to pick one
of them or ``tip something being carried''.
\\
%.lp ""
If the latter is chosen, there will be another prompt for which item
from inventory to tip.
The `\texttt{m}' prefix makes the command skip containers on the
floor and pick one from inventory, except for the special case of
\textit{menustyle:Traditional}
with two or more containers present; that situation will start with the
floor container menu.
\\
%.lp ""
Autocompletes. Default key is `\texttt{M-T}'.
%.lp
\item[\#toggle]
Toggle a boolean option on or off.
Requires a parameter in parenthesis, the name of the option to toggle.
The option must be settable in-game.

%.lp ""
For example:
%.sd
\begin{verbatim}
    BIND=':toggle(price_quotes)
    BIND=@:toggle(autopickup)
\end{verbatim}
%.ed

%.lp
\item[\#travel]
Travel to a specific location on the map.
Default key is `\texttt{\textunderscore }'.
Using the ``request menu'' prefix shows a menu of interesting targets in sight
without asking to move the cursor.
When picking a target with cursor and the \textit{autodescribe}
option is on, the top line will show ``(no travel path)'' if
your character does not know of a path to that location.
See also \texttt{\#retravel}.
%.lp
\item[\#turn]
Turn undead away. Autocompletes. Default key is `\texttt{M-t}'.
%.lp
\item[\#twoweapon]
Toggle two-weapon combat on or off. Autocompletes.
Default key is `\texttt{X}',
and also `\texttt{M-2}' if \textit{number\textunderscore pad} is off.\\
%.lp ""
Note that you must
use suitable weapons for this type of combat, or it will
be automatically turned off.
%.lp
\item[\#untrap]
Untrap something (trap, door, or chest).
Default key is `\texttt{M-u}', and `\texttt{u}' if \textit{number\textunderscore pad} is on.\\
%.lp ""
In some circumstances it can also be used to rescue trapped monsters.
%.lp
\item[\#up]
Go up a staircase. Default key is `\texttt{<}'.
%.lp
\item[\#vanquished]
List vanquished monsters by type and count.
\\
%.lp ""
Note that the vanquished monsters list includes all monsters killed by
traps and each other as well as by you, and omits any which got removed
from the game without being killed (perhaps by genocide, or by a mollified
shopkeeper dismissing summoned Kops) or were already corpses when placed
on the map.
\\
%.lp ""
Using the ``request menu'' prefix prior to \#vanquished brings up
a menu of sorting orders available (provided that the vanquished monsters
list contains at least two types of monsters).
Whichever ordering is picked gets assigned to the \textit{sortvanquished}
option so is remembered for subsequent \#vanquished requests.
The \texttt{\#genocided} command shares this sorting order.
\\
%.lp ""
During end-of-game disclosure, when asked whether to show vanquished
monsters answering `\texttt{a}' will let you choose from the sort menu.
\\
%.lp ""
Autocompletes.
Default key is `\texttt{M-V}'.
%.lp
\item[\#version]
Print compile time options for this version of \textit{NetHack}.

%.lp
The second paragraph lists the user interface(s) that are included.
If there are more than one, you can use the \textit{windowtype}
option in your run-time configuration file to select the one you want.

%.lp
Autocompletes. Default key is `\texttt{M-v}'.
%.lp
\item[\#versionshort]
Show the program's version number, plus the date and time that the
running copy was built from sources (not the version's release date).
Default key is `\texttt{V}'.
%.lp
\item[\#vision]
Show vision array.
Autocompletes.
Debug mode only.
%.lp
\item[\#wait]
Rest one move while doing nothing.
Default key is `\texttt{.}', and also `{\tt{ }}' if
\textit{rest\textunderscore on\textunderscore space} is on.
%.lp
\item[\#wear]
Wear a piece of armor. Default key is `\texttt{W}'.
%.lp
\item[\#whatdoes]
Tell what a key does. Default key is `\texttt{\&}'.
%.lp
\item[\#whatis]
Show what type of thing a symbol corresponds to. Default key is `\texttt{/}'.
%.lp
\item[\#wield]
Wield a weapon. Default key is `\texttt{w}'.
%.lp
\item[\#wipe]
Wipe off your face. Autocompletes. Default key is `\texttt{M-w}'.
%.lp
\item[\#wizborn]
Show monster birth, death, genocide, and extinct statistics.
Debug mode only.
%.lp
\item[\#wizbury]
Bury objects under and around you.
Autocompletes.
Debug mode only.
%.lp
\item[\#wizcast]
Cast any spell.
Debug mode only.
%.lp
\item[\#wizdetect]
Reveal hidden things (secret doors or traps or unseen monsters)
within a modest radius.
No time elapses.
Autocompletes.
Debug mode only.
Default key is `\texttt{\textasciicircum E}'.
%.lp
\item[\#wizgenesis]
Create a monster.
May be prefixed by a count to create more than one.
Autocompletes.
Debug mode only.
Default key is `\texttt{\textasciicircum G}'.
%.lp
\item[\#wizidentify]
Identify all items in inventory.
Autocompletes.
Debug mode only.
Default key is `\texttt{\textasciicircum I}'.
%.lp
\item[\#wizintrinsic]
Set one or more intrinsic attributes.
Autocompletes.
Debug mode only.
%.lp
\item[\#wizkill]
Remove monsters from play by just pointing at them.
By default the hero gets credit or blame for killing the targets.
Precede this command with the `\texttt{m}' prefix to override that.
Autocompletes.
Debug mode only.
%.lp
\item[\#wizlevelport]
Teleport to another level.
Autocompletes.
Debug mode only.
Default key is `\texttt{\textasciicircum V}'.
%.lp
\item[\#wizmap]
Map the level.
Autocompletes.
Debug mode only.
Default key is `\texttt{\textasciicircum F}'.
%.lp
\item[\#wizrumorcheck]
Verify rumor boundaries by displaying first and last true rumors and
first and last false rumors.\\
%.lp ""
Also displays first, second, and last random engravings, epitaphs,
and hallucinatory monsters.\\
%.lp ""
Autocompletes.
Debug mode only.
%.lp
\item[\#wizseenv]
Show map locations' seen vectors.
Autocompletes.
Debug mode only.
%.lp
\item[\#wizsmell]
Smell monster.
Autocompletes.
Debug mode only.
%.lp
\item[\#wizwhere]
Show locations of special levels.
Autocompletes.
Debug mode only.
%.lp
\item[\#wizwish]
Wish for something.
Autocompletes.
Debug mode only.
Default key is `\texttt{\textasciicircum W}'.
Precede this command with the `\texttt{m}' prefix to show a wish history menu.
%.lp
\item[\#wmode]
Show wall modes.
Autocompletes.
Debug mode only.
%.lp
\item[\#zap]
Zap a wand. Default key is `\texttt{z}'.
%.lp
\item[\#?]
Help menu:  get the list of available extended commands.
\end{description}

%.pg
\noindent If your keyboard has a meta key (which, when pressed in combination
with another key, modifies it by setting the `meta' [8th, or `high']
bit), you can invoke many extended commands by meta-ing the first
letter of the command.

On \textit{Windows} and \textit{MS-DOS},
the `Alt' key can be used in this fashion.
On other systems, if typing `Alt' plus another key transmits a
two character sequence consisting of an \texttt{Escape}
followed by the other key, you may set the \textit{altmeta}
option to have \textit{NetHack} combine them into \texttt{meta+<key>}.
(This combining action only takes place when NetHack is expecting a
command to execute, not when accepting input to name something or to
make a wish.)

%.pg
Unlike control characters, where \texttt{\textasciicircum x} and \texttt{\textasciicircum X} denote the same
thing, meta characters are case-sensitive:  \texttt{M-x} and \texttt{M-X}
represent different things.  Some commands which can be run via a meta
character require that the letter be capitalized because the lower-case
equivalent is used for another command, so the three key combination
\texttt{meta+Shift+letter} is needed.

%.BR 1
\begin{description}[font=\mdseries\ttfamily]
%.lp
\item[M-?]
{\tt\#?} (not supported by all platforms)
%.lp
\item[M-2]
{\tt\#twoweapon} (unless the \textit{number\textunderscore pad} option is enabled)
%.lp
\item[M-a]
{\tt\#adjust}
%.lp
\item[M-A]
{\tt\#annotate}
%.lp
\item[M-c]
{\tt\#chat}
%.lp
\item[M-C]
{\tt\#conduct}
%.lp
\item[M-d]
{\tt\#dip}
%.lp
\item[M-e]
{\tt\#enhance}
%.lp
\item[M-f]
{\tt\#force}
%.lp
\item[M-g]
{\tt\#genocided}
%.lp
\item[M-i]
{\tt\#invoke}
%.lp
\item[M-j]
{\tt\#jump}
%.lp
\item[M-l]
{\tt\#loot}
%.lp
\item[M-m]
{\tt\#monster}
%.lp
\item[M-n]
{\tt\#name}
%.lp
\item[M-o]
{\tt\#offer}
%.lp
\item[M-O]
{\tt\#overview}
%.lp
\item[M-p]
{\tt\#pray}
%.Ip
\item[M-r]
{\tt\#rub}
%.lp
\item[M-R]
{\tt\#ride}
%.lp
\item[M-s]
{\tt\#sit}
%.lp
\item[M-t]
{\tt\#turn}
%.lp
\item[M-T]
{\tt\#tip}
%.lp
\item[M-u]
{\tt\#untrap}
%.lp
\item[M-v]
{\tt\#version}
%.lp
\item[M-V]
{\tt\#vanquished}
%.lp
\item[M-w]
{\tt\#wipe}
%.lp
\item[M-X]
{\tt\#exploremode}
\end{description}

%.pg
\noindent If the \textit{number\textunderscore pad} option is on, some additional letter commands
are available:
\begin{description}[font=\mdseries\ttfamily]
%.lp
\item[h]
{\tt\#help}
%.lp
\item[j]
{\tt\#jump}
%.lp
\item[k]
{\tt\#kick}
%.lp
\item[l]
{\tt\#loot}
%.lp
\item[N]
{\tt\#name}
%.lp
\item[u]
{\tt\#untrap}
\end{description}

%.BR 1 \"blank line for extra separation; plain text output looks better

%.hn 1
\section{Rooms and corridors}

%.pg
Rooms and corridors in the dungeon are either lit or dark.
Any lit areas within your line of sight will be displayed;
dark areas are only displayed if they are within one space of you.
Walls and corridors remain on the map as you explore them.

%.pg
Secret corridors are hidden and appear to be solid rock.
You can find them with the `\texttt{s}' (search) command when adjacent
to them.
Multiple search attempts may be needed.
When searching is successful, secret corridors become ordinary open
corridor locations.
Mapping magic reveals secret corridors, so converts them into ordinary
corridors and shows them as such.

%.hn 2
\subsection*{Doorways}

%.pg
Doorways connect rooms and corridors.
Some doorways have no doors; you can walk right through.
Others have doors in them, which may be open, closed, or locked.
To open a closed door, use the `\texttt{o}' (open)
command; to close it again, use the `\texttt{c}' (close) command.
By default the
\textit{autoopen}
option is enabled, so simply attempting to walk onto a closed door's
location will attempt to open it without needing `\texttt{o}'.
Opening via
\textit{autoopen}
will not work if you are \textit{confused} or \textit{stunned} or suffer from
the \textit{fumbling} attribute.

%.pg
Open doors cannot be entered diagonally; you must approach them
straight on, horizontally or vertically.
Doorways without doors are
not restricted in this fashion except on one particular level
%.\" the rogue level
(described by ``\texttt{\#overview}'' as ``a primitive area'').

%.pg
Unlocking magic exists but usually won't be available early on.
You can get through a locked door without magic by first using an
unlocking tool with the `\texttt{a}' (apply) command, and then opening it.
By default the
\textit{autounlock}
option is also enabled, so if you attempt to open (via `\texttt{o}' or
\textit{autoopen})
a locked door while carrying an unlocking tool, you'll be asked whether
to use it on the door's lock.
Alternatively, you can break a closed door (whether locked or not) down
by kicking it via the ``\texttt{\textasciicircum D}'' (kick) command.
Kicking down a door destroys it and makes a lot of noise which might
wake sleeping monsters.

%.pg
Some closed doors are booby-trapped and will explode if an attempt is made
to open (when unlocked) or unlock (when locked) or kick down.
Like kicking, an explosion destroys the door and makes a lot of noise.
The ``\texttt{\#untrap}'' command can be used to search a door for traps but
might take multiple attempts to find one.
When one is found, you'll be asked whether to try to disarm it.
If you accede, success will eliminate the trap but
failure will set off the trap's explosion.
(If you decline, you effectively forget that a trap was found there.)

%.pg
Closed doors can be useful for shutting out monsters.
Most monsters cannot open closed doors, although a few don't need to
(for example, ghosts can walk through doors and fog clouds can flow
under them).
Some monsters who can open doors can also use unlocking tools.
And some (giants) can smash doors.

%.pg
Secret doors are hidden and appear to be ordinary wall (from inside a
room) or solid rock (from outside).
You can find them with the `\texttt{s}' (search) command but it might
take multiple tries (possibly many tries if your luck is poor).
Once found they are in all ways equivalent to normal doors.
Mapping magic does not reveal secret doors.

%.hn 2
\subsection*{Traps (`\texttt{\textasciicircum }')}

%.pg
There are traps throughout the dungeon to snare the unwary intruder.
For example, you may suddenly fall into a pit and be stuck for a few
turns trying to climb out (see below).
A trap usually won't appear on your map until you trigger it by moving
onto it, you see someone else trigger it, or you discover it with
the `\texttt{s}' (search) command (multiple attempts are often needed;
if your luck is poor, many attempts might be needed).
\textit{Wands of secret door detection} and the spell of \textit{detect unseen}
also reveal traps within a modest radius but only if the trap is also within
line-of-sight (whether you can see at the time or not).
There is also other magic which can reveal traps.

%.pg
Monsters can fall prey to
traps, too, which can potentially be used as a defensive strategy.
Unfortunately traps can be harmful to your pet(s) as well.
Monsters, including pets, usually will avoid moving onto a trap which
is shown on your map if they have encountered that type of trap before.

%.pg
Some traps such as pits, bear traps, and webs hold you in one place.
You can escape by simply trying to move to an adjacent spot and repeat
as needed; eventually you will get free.

%.pg
Other traps can send you to different locations.
Teleporters send you elsewhere on the same dungeon level.
Level teleporters send you to a random dungeon level, the destination
chosen from a few levels lower all the way to the top.
These traps choose a new destination each time they're activated.
Trap doors and holes also send you to another level, but one which is
always below the current level.
Usually that will be the next level down but it can be farther.
Unlike (level) teleporters, the destination level of a particular trap door
or hole is persistent, so falling into one will bring you to the same level
each time---though not necessarily the same spot on the level.
Magic portals behave similarly, but with some additional variation.
Some portals are two-way and their remote destination is always the same:
another portal which can take you back.
Others are one-way and send you to a specific destination level but not
necessarily to a specific location there.

%.pg
There is a special multi-level branch of the dungeon with pre-mapped levels
based on the classic computer game ``\textit{Sokoban}.''
In that game, you operate as a warehouse worker who pushes crates around
obstacles to position them at designated locations.
In NetHack, the goal is to push boulders into pits or holes until those
traps have all been nullified, giving access to whatever is beyond them.
In the Sokoban game, you can only move in the four cardinal compass
directions, and a crate in its final destination blocks further access
to that spot.
In the Sokoban levels of NetHack, you can move diagonally (unless that
would let you pass between two neighboring boulders) but you can only
push boulders in the four cardinal directions, and a boulder which fills
a pit or hole removes both the boulder and the trap so opens up normal
access to that spot.
With careful foresight, it is possible to complete all of the levels
according to the traditional rules of Sokoban.
(Hint: to solve Sokoban puzzles, you often need to move things away from
their eventual destinations in order to open up more room to maneuver.)
Since NetHack does not support an \textit{undo} capability, some allowances
are permitted in case you get stuck.
For example, each level has at least one extra boulder.
Also, it is possible to drop everything in order to be able to squeeze
into the same location as a boulder (and then presumably move past it),
or to destroy a boulder with magic or tools, or to create new boulders
with a \textit{scroll of earth}.
However, doing such things will lower your luck without any specific
message given about that.
See the \textit{Conduct} section for information about getting feedback for
your actions in Sokoban.

%.hn 2
\subsection*{Stairs and ladders (`\texttt{<}', `\texttt{>}')}

%.pg
In general, each level in the dungeon will have a staircase going up
(`\texttt{<}') to the previous level and another going down (`\texttt{>}')
to the next
level.  There are some exceptions though.  For instance, fairly early
in the dungeon you will find a level with two down staircases, one
continuing into the dungeon and the other branching into an area
known as the Gnomish Mines.  Those mines eventually hit a dead end,
so after exploring them (if you choose to do so), you'll need to
climb back up to the main dungeon.

%.pg
When you traverse a set of stairs, or trigger a trap which sends you
to another level, the level you're leaving will be deactivated and
stored in a file on disk.  If you're moving to a previously visited
level, it will be loaded from its file on disk and reactivated.  If
you're moving to a level which has not yet been visited, it will be
created (from scratch for most random levels, from a template for
some ``special'' levels, or loaded from the remains of an earlier game
for a ``bones'' level as briefly described below).  Monsters are only
active on the current level; those on other levels are essentially
placed into stasis.

%.pg
Ordinarily when you climb a set of stairs, you will arrive on the
corresponding staircase at your destination.  However, pets (see below)
and some other monsters will follow along if they're close enough when
you travel up or down stairs, and occasionally one of these creatures
will displace you during the climb.  When that occurs, the pet or other
monster will arrive on the staircase and you will end up nearby.

%.pg
Ladders serve the same purpose as staircases, and the two types of
inter-level connections are nearly indistinguishable during game play.

%.hn 2
\subsection*{Shops and shopping}

%.pg
Occasionally you will run across a room with a shopkeeper near the door
and many items lying on the floor.  You can buy items by picking them
up and then using the `\texttt{p}' command.  You can inquire about the price
of an item prior to picking it up by using the ``\texttt{\#chat}'' command
while standing on it.  Using an item prior to paying for it will incur a
charge, and the shopkeeper won't allow you to leave the shop until you
have paid any debt you owe.

%.pg
You can sell items to a shopkeeper by dropping them to the floor while
inside a shop.  You will either be offered an amount of gold and asked
whether you're willing to sell, or you'll be told that the shopkeeper
isn't interested (generally, your item needs to be compatible with the
type of merchandise carried by the shop).

%.pg
If you drop something in a shop by accident, the shopkeeper will usually
claim ownership without offering any compensation.  You'll have to buy
it back if you want to reclaim it.

%.pg
Shopkeepers sometime run out of money.
When that happens, you'll be
offered credit instead of gold when you try to sell something.
Credit can be used to pay for purchases, but it is only good in the shop
where it was obtained; other shopkeepers won't honor it.
(If you happen to
find a ``credit card'' in the dungeon, don't bother trying to use it in
shops; shopkeepers will not accept it.)

%.pg
The \texttt{\$} command, which reports the amount of gold you are carrying,
will also show current shop debt or credit, if any.
The \texttt{Iu} command lists unpaid items (those which still belong to the
shop) if you are carrying any.
The \texttt{Ix} command shows an inventory-like display of any unpaid items
which have been used up, along with other shop fees, if any.

%.hn 3
\subsubsection*{Shop idiosyncrasies}

%.pg
Several aspects of shop behavior might be unexpected.

\begin{itemize}
% note: a bullet is the default item label so we could omit [$\bullet$] here
%.lp \(bu 2
\item
The price of a given item can vary due to a variety of factors.
%.lp \(bu 2
\item
A shopkeeper treats the spot immediately inside the door as if it were
outside the shop.
%.lp \(bu 2
\item
While the shopkeeper watches you like a hawk, he or she will generally ignore
any other customers.
%.lp \(bu 2
\item
If a shop is ``closed for inventory,'' it will not open of its own accord.
%.lp \(bu 2
\item
Shops do not get restocked with new items, regardless of inventory depletion.
\end{itemize}

%.hn 2
\subsection*{Movement feedback}

%.pg
Moving around the map usually provides no
feedback---other than drawing the hero at the new location---unless
you step on an object or pile of objects,
or on a trap, or attempt to move onto a spot where a monster is located.
There are several options which can be used to augment the normal feedback.

%.pg
The
\textit{pile\textunderscore limit}
option controls how many objects can be in a
pile---sharing the same map location---for
the game to state ``there are objects here'' instead of listing them.
The default is \texttt{5}.
Setting it to \texttt{1} would always give that message instead of listing
any objects.
Setting it to \texttt{0} is a special case which will always list all
objects no matter how big a pile is.
Note that the number refers to the count of separate stacks of objects
present rather than the sum of the quantities of those stacks (so
\texttt{7 arrows} or \texttt{25 gold pieces} will each count as 1 rather
than as 7 and 25, respectively, and total to 2 when both are at the
same location).

%.pg
The \texttt{nopickup} command prefix (default `\texttt{m}') can be
used before a movement direction to step on objects without attempting
auto-pickup and without giving feedback about them.

%.pg
The
\textit{mention\textunderscore walls}
option controls whether you get feedback if you try to walk into a wall
or solid stone or off the edge of the map.
Normally nothing happens (unless the hero is blind and no wall is shown,
then the wall that is being bumped into will be drawn on the map).
This option also gives feedback when rushing or running stops for
some non-obvious reason.

%.pg
The
\textit{mention\textunderscore decor}
option controls whether you get feedback when walking on ``furniture.''
Normally stepping onto stairs or a fountain or an altar or various other
things doesn't elicit anything unless it is covered by one or more objects
so is obscured on the map.
Setting this option to true will describe such things even when they
aren't obscured.
Doorless doorways and open doors aren't considered worthy of mention;
closed doors (if you can move onto their spots) and broken doors are.
Assuming that you're able to do so, moving onto water or lava or ice
will give feedback if not yet on that type of terrain but not repeat it
(unless there has been some intervening message) when moving from water
to another water spot, or lava to lava, or ice to ice.
Moving off of any of those back onto ``normal'' terrain will give one
message too, unless there is feedback about one or more objects, in which
case the back on land circumstance is implied.

%.pg
The
\textit{confirm}
and
\textit{safe\textunderscore pet}
options control what happens when you try to move onto a peaceful monster's
spot or a tame one's spot.

%.\" getting away from "Movement feedback" here; oh well...
%.pg
The \texttt{nopickup} command prefix (default `\texttt{m}') is
also the move-without-attacking prefix and can be used to try to step
onto a visible monster's spot without the move being considered an attack
(see the \textit{Fighting} subsection of \textit{Monsters} below).
The `\texttt{fight}' command prefix (default `\texttt{F}';
also `\texttt{-}' if
\textit{number\textunderscore pad}
is on) can be used to force an attack, when guessing where an unseen
monster is or when deliberately attacking a peaceful or tame creature.

%.pg
The
\textit{run\textunderscore mode}
option controls how frequently the map gets redrawn when moving more
than one step in a single command (so when rushing, running, or traveling).

%.hn 2
\subsection*{Rogue level}

%.pg
One dungeon level (occurring in mid to late teens of the main dungeon)
is a tribute to the ancestor game \textit{hack}'s inspiration \textit{rogue}.

%.pg
It is usually displayed differently from other levels: possibly in
characters instead of tiles, or without line-drawing symbols if already
in characters; also, gold is shown as \texttt{*} rather than \texttt{\textdollar}
and stairs are shown as \texttt{\%} rather than \texttt{<} and \texttt{>}.
There are some minor differences in actual game play: doorways lack
doors; a scroll, wand, or spell of light used in a room lights up the
whole room rather than within a radius around your character.
And monsters represented by lower-case letters aren't randomly
generated on the rogue level.

%.pg
The slight strangeness of this level is a feature, not a bug....

%.hn 1
\section{Monsters}

%.pg
Monsters you cannot see are not displayed on the screen.  Beware!
You may suddenly come upon one in a dark place.  Some magic items can
help you locate them before they locate you (which some monsters can do
very well).

%.pg
The commands `\texttt{/}' and `\texttt{;}' may be used to obtain information
about those
monsters who are displayed on the screen.  The command ``\texttt{\#name}''
(by default bound to `\texttt{C}'), allows you
to assign a name to a monster, which may be useful to help distinguish
one from another when multiple monsters are present.  Assigning a name
which is just a space will remove any prior name.

%.pg
The extended command ``\texttt{\#chat}'' can be used to interact with an adjacent
monster.  There is no actual dialog (in other words, you don't get to
choose what you'll say), but chatting with some monsters such as a
shopkeeper or the Oracle of Delphi can produce useful results.

%.hn 2
\subsection*{Fighting}

%.pg
If you see a monster and you wish to fight it, just attempt to walk
into it.  Many monsters you find will mind their own business unless
you attack them.  Some of them are very dangerous when angered.
Remember:  discretion is the better part of valor.

%.pg
In most circumstances, if you attempt to attack a peaceful monster by
moving into its location, you'll be asked to confirm your intent.  By
default an answer of `\texttt{y}' acknowledges that intent,
which can be error prone if you're using `\texttt{y}' to move.  You can set the
\textit{paranoid\textunderscore confirmation:attack}
option to require a response of ``\texttt{yes}'' instead.
%.pg

If you can't see a monster (if it is invisible, or if you are blinded),
the symbol `\texttt{I}' will be shown when you learn of its presence.
If you attempt to walk into it, you will try to fight it just like
a monster that you can see; of course,
if the monster has moved, you will attack empty air.  If you guess
that the monster has moved and you don't wish to fight, you can use the
`\texttt{m}' command to move without fighting; likewise, if you don't remember
a monster but want to try fighting anyway, you can use the `\texttt{F}' command.

%.hn 2
\subsection*{Your pet}

%.pg
You start the game with a little dog (`\texttt{d}'), kitten (`\texttt{f}'),
or pony (`\texttt{u}'), which follows
you about the dungeon and fights monsters with you.
Like you, your pet needs food to survive.
Dogs and cats usually feed themselves on fresh carrion and other meats;
horses need vegetarian food which is harder to come by.
If you're worried about your pet or want to train it, you
can feed it, too, by throwing it food.
A properly trained pet can be very useful under certain circumstances.

%.pg
Your pet also gains experience from killing monsters, and can grow
over time, gaining hit points and doing more damage.  Initially, your
pet may even be better at killing things than you, which makes pets
useful for low-level characters.

%.pg
Your pet will follow you up and down staircases if it is next to you
when you move.  Otherwise your pet will be stranded and may become
wild.  Similarly, when you trigger certain types of traps which alter
your location (for instance, a trap door which drops you to a lower
dungeon level), any adjacent pet will accompany you and any non-adjacent
pet will be left behind.  Your pet may trigger such traps itself; you
will not be carried along with it even if adjacent at the time.

%.hn 2
\subsection*{Steeds}

%.pg
Some types of creatures in the dungeon can actually be ridden if you
have the right equipment and skill.  Convincing a wild beast to let
you saddle it up is difficult to say the least.  Many a dungeoneer
has had to resort to magic and wizardry in order to forge the alliance.
Once you do have the beast under your control however, you can
easily climb in and out of the saddle with the ``\texttt{\#ride}'' command.  Lead
the beast around the dungeon when riding, in the same manner as
you would move yourself.  It is the beast that you will see displayed
on the map.

%.pg
Riding skill is managed by the ``\texttt{\#enhance}'' command.  See the section
on Weapon proficiency for more information about that.

%.pg
Use the `\texttt{a}' (apply) command and pick a saddle in your inventory to
attempt to put that saddle on an adjacent creature.  If successful,
it will be transferred to that creature's inventory.

%.pg
Use the ``\texttt{\#loot}'' command while adjacent to a saddled creature to
try to remove the saddle from that creature.  If successful, it will
be transferred to your inventory.

%.hn 2
\subsection*{Bones levels}

%.pg
You may encounter the shades and corpses of other adventurers (or even
former incarnations of yourself!) and their personal effects.  Ghosts
are hard to kill, but easy to avoid, since they're slow and do little
damage.  You can plunder the deceased adventurer's possessions;
however, they are likely to be cursed.  Beware of whatever killed the
former player; it is probably still lurking around, gloating over its
last victory.

%.hn 2
\subsection*{Persistence of Monsters}

%.pg
Monsters (a generic reference which also includes humans and pets) are only
shown while they can be seen or otherwise sensed.
Moving to a location where you can't see or sense a monster any more
will result in it disappearing from your map, similarly if it is the
one who moved rather than you.

%.pg
However, if you encounter a monster which you can't see or
sense---perhaps it is invisible and has just tapped you on the noggin---a
special ``remembered, unseen monster'' marker will be displayed at
the location where you think it is.
That will persist until you have
proven that there is no monster there, even if the unseen monster
moves to another location or you move to a spot where the marker's
location ordinarily wouldn't be seen any more.

%.hn 1
\section{Objects}

%.pg
When you find something in the dungeon, it is common to want to pick
it up.  In \textit{NetHack}, this is accomplished by using the `\texttt{,}' command.
If \textit{autopickup} option is on, you will automatically pick up the object
by walking over it, unless you move with the `\texttt{m}' prefix.
%.pg
If you're carrying too many items, \textit{NetHack} will tell you so and you
won't be able to pick up anything more.  Otherwise, it will add the object(s)
to your pack and tell you what you just picked up.
%.pg
As you add items to your inventory, you also add the weight of that object
to your load.  The amount that you can carry depends on your strength and
your constitution.  The
stronger and sturdier
you are, the less the additional load will affect you.  There comes
a point, though, when the weight of all of that stuff you are carrying around
with you through the dungeon will encumber you.  Your reactions
will get slower and you'll burn calories faster, requiring food more frequently
to cope with it.  Eventually, you'll be so overloaded that you'll either have
to discard some of what you're carrying or collapse under its weight.
%.pg
\textit{NetHack} will tell you how badly you have loaded yourself.
If you are encumbered, one of the conditions
\texttt{Burdened}, \texttt{Stressed}, \texttt{Strained},
\texttt{Overtaxed}, or \texttt{Overloaded} will be
shown on the bottom line status display.

%.pg
When you pick up an object, it is assigned an inventory letter.  Many
commands that operate on objects must ask you to find out which object
you want to use.  When \textit{NetHack} asks you to choose a particular object
you are carrying, you are usually presented with a list of inventory
letters to choose from (see Commands, above).

%.pg
Some objects, such as weapons, are easily differentiated.  Others, like
scrolls and potions, are given descriptions which vary according to
type.  During a game, any two objects with the same description are
the same type.  However, the descriptions will vary from game to game.

%.pg
When you use one of these objects, if its effect is obvious, \textit{NetHack}
will remember what it is for you.  If its effect isn't extremely
obvious, you will be asked what you want to call this type of object
so you will recognize it later.  You can also use the ``\texttt{\#name}''
command, for the same purpose at any time, to name
all objects of a particular type or just an individual object.
When you use ``\texttt{\#name}'' on an object which has already been named,
specifying a space as the value will remove the prior name instead
of assigning a new one.

%.hn 2
\subsection*{Curses and Blessings}

%.pg
Any object that you find may be cursed, even if the object is
otherwise helpful.  The most common effect of a curse is being stuck
with (and to) the item.  Cursed weapons weld themselves to your hand
when wielded, so you cannot unwield them.  Any cursed item you wear
is not removable by ordinary means.  In addition, cursed arms and armor
usually, but not always, bear negative enchantments that make them
less effective in combat.  Other cursed objects may act poorly or
detrimentally in other ways.

%.pg
Objects can also be blessed instead.
Blessed items usually work better or more beneficially than normal
uncursed items.
For example, a blessed weapon will do slightly more damage against demons.

%.pg
Objects which are neither cursed nor blessed are referred to as uncursed.
They could just as easily have been described as unblessed, but the
uncursed designation is what you will see within the game.  A ``glass
half full versus glass half empty'' situation; make of that what you will.

%.pg
There are magical means of bestowing or removing curses upon objects,
so even if you are stuck with one, you can still have the curse
lifted and the item removed.
Priests and Priestesses have an innate
sensitivity to this property in any object, so they can more easily avoid
cursed objects than other character roles.
Dropping objects onto an altar will reveal their bless or curse state
provided that you can see them land.

%.pg
An item with unknown status will be reported in your inventory with no prefix.
An item which you know the state of will be distinguished in your inventory
by the presence of the word \texttt{cursed}, \texttt{uncursed} or
\texttt{blessed} in the description of the item.
In some cases \texttt{uncursed} will be omitted as being redundant when
enough other information is displayed.
The
\textit{implicit\textunderscore uncursed}
option can be used to control this; toggle it off to have \texttt{uncursed}
be displayed even when that can be deduced from other attributes.

%.pg
Sometimes the bless or curse state of objects is referred to as their
``\texttt{BUC}'' attribute, for Blessed, Uncursed, or Cursed state,
or ``\texttt{BUCX}'' for Blessed, Uncursed, Cursed, or unknown.
(The term \textit{beatitude} is occasionally used as well.)

%.hn 2
\subsection*{Artifacts}

%.pg
Some objects have been imbued with special powers and are known as
\textit{Artifacts}.
They have specific types (such as long sword or orcish dagger) and distinct
names such as \textit{Giantslayer} or \textit{Grimtooth}.
Artifact weapons typically do more damage than their ordinary counterparts.
Some do extra damage against all monsters, others only against specific
types of monsters so aren't better than regular weapons against other types.
Some confer defensive capabilities when wielded or have other powers that
aren't listed here.

%.pg
You might find them simply lying on the floor, including but not limited
to inside shops, or be granted as a reward for ``\texttt{\#offer}'' on an
altar to your patron deity.
A few might be dropped by monsters, or might be converted from an ordinary
object of the same type via assigning the right name.
% should we mention dipping for Excalibur here?
Or you might wish for them, if you happen to be granted a wish, but such
wishes can fail.

%.pg
Some artifacts have a specific alignment, others don't.
You won't obtain aligned ones that have a different alignment from yours
via offering and might get a shock if you attempt to wish for any of those
or find one and attempt to use it.

%pg
Each role has a distinct artifact that is contained in the \textit{Quest}
dungeon branch.
These are commonly known as quest artifacts.
All are aligned and most are non-weapons.
They won't be found randomly.

%.pg
The `\texttt{\textbackslash}' and `\texttt{\textasciigrave{}a}' commands will
list artifacts that you have fully identified (knowing the name and item
type isn't sufficient).

%.hn 2
\subsection*{Relics}

%.pg
There are three unique items that are named and have limited special
powers but aren't classified as artifacts.
Each is guarded by a particular monster and you'll need to collect all
three for use late in the game.
% The relics are listed in the same order as the Oracle's message about
% them rather than in the order they need to be used for the invocation.
They are \textit{the Bell of Opening},
\textit{the Book of the Dead}, and
\textit{the Candelabrum of Invocation}.
Their corresponding descriptions when not yet identified are
silver bell, papyrus spellbook, and candelabrum.

%.hn 2
\subsection*{Weapons (`\texttt{)}')}

%.pg
Given a chance, most monsters in the Mazes of Menace will gratuitously try to
kill you.
You need weapons for self-defense (killing them first).
Without a
weapon, you do only 1--2 hit points of damage (plus bonuses, if any).
Monk characters are an exception; they normally do more damage with
bare (or gloved) hands than they do with weapons.

%.pg
There are wielded weapons, like maces and swords, and thrown weapons,
like arrows and spears.  To hit monsters with a weapon, you must wield it and
attack them, or throw it at them.  You can simply elect to throw a spear.
To shoot an arrow, you should first wield a bow, then throw the arrow.
Crossbows shoot crossbow bolts.  Slings hurl rocks and (other) stones
(like gems).

%.pg
Enchanted weapons have a ``plus'' (or ``to hit enhancement'' which can be
either positive or negative) that adds to your chance to
hit and the damage you do to a monster.  The only way to determine a weapon's
enchantment is to have it magically identified somehow.
Most weapons are subject to some type of damage like rust.  Such
``erosion'' damage can be repaired.

%.pg
The chance that an attack will successfully hit a monster, and the amount
of damage such a hit will do, depends upon many factors.  Among them are:
type of weapon, quality of weapon (enchantment and/or erosion), experience
level, strength, dexterity, encumbrance, and proficiency (see below).  The
monster's armor
class---a general defense rating, not necessarily due to wearing of armor---is
a factor too; also, some monsters are particularly
vulnerable to certain types of weapons.

%.pg
Many weapons can be wielded in one hand; some require both hands.
When wielding a two-handed weapon, you can not wear a shield, and
vice versa.  When wielding a one-handed weapon, you can have another
weapon ready to use by setting things up with the `\texttt{x}' command, which
exchanges your primary (the one being wielded) and alternate weapons.
And if you have proficiency in the ``two weapon combat'' skill, you
may wield both weapons simultaneously as primary and secondary; use the
`\texttt{X}' command to engage or disengage that.
Only some types of characters (barbarians, for instance) have the necessary
skill available.  Even with that skill, using two weapons at once incurs
a penalty in the chance to hit your target compared to using just one
weapon at a time.

%.pg
There might be times when you'd rather not wield any weapon at all.
To accomplish that, wield `\texttt{-}', or else use the `\texttt{A}' command which
allows you to unwield the current weapon in addition to taking off
other worn items.

%.pg
Those of you in the audience who are AD\&D players, be aware that each
weapon which existed in AD\&D does roughly the same damage to monsters in
\textit{NetHack}.  Some of the more obscure weapons (such as the
\textit{aklys}, \textit{lucern hammer}, and \textit{bec-de-corbin}) are defined
in an appendix to \textit{Unearthed Arcana}, an AD\&D supplement.

%.pg
Some interfaces support the \texttt{weaponstatus} option.
When it is enabled, an extra status condition is displayed, describing
the currently wielded weapon.

%.pg
The commands to use weapons are `\texttt{w}' (wield), `\texttt{t}' (throw),
`\texttt{f}' (fire), `\texttt{Q}' (quiver),
`\texttt{x}' (exchange), `\texttt{X}' (twoweapon), and ``\texttt{\#enhance}''
(see below).

%.hn 3
\subsection*{Throwing and shooting}

%.pg
You can throw just about anything via the `\texttt{t}' command.  It will prompt
for the item to throw; picking `\texttt{?}' will list things in your inventory
which are considered likely to be thrown, or picking `\texttt{*}' will list
your entire inventory.  After you've chosen what to throw, you will
be prompted for a direction rather than for a specific target.  The
distance something can be thrown depends mainly on the type of object
and your strength.  Arrows can be thrown by hand, but can be thrown
much farther and will be more likely to hit when thrown while you are
wielding a bow.

%.pg
Some weapons will return when thrown.
A boomerang---provided it fails to hit anything---is an obvious example.
If an aklys (thonged club) is thrown while it is wielded, it will return
even when it hits something.
A sufficiently strong hero can throw the warhammer \textit{Mjollnir};
when thrown by a \textit{Valkyrie} it will return too.
However, aklyses and \textit{Mjollnir} occasionally fail to return.
Returning thrown objects occasionally fail to be caught, sometimes even
hitting the thrower, but when caught they become re-wielded.

%.pg
You can simplify the throwing operation by using the `\texttt{Q}' command to
select your preferred ``missile'', then using the `\texttt{f}' command to
throw it.  You'll be prompted for a direction as above, but you don't
have to specify which item to throw each time you use `\texttt{f}'.  There is
also an option,
\textit{autoquiver},
which has \textit{NetHack} choose another item to automatically fill your
quiver (or quiver sack, or have at the ready) when the inventory slot used
for `\texttt{Q}' runs out.
If your quiver is empty, \textit{autoquiver}
is false, and you are wielding a weapon which returns when thrown,
you will throw that weapon instead of filling the quiver.
The fire command also has extra assistance, if \textit{fireassist}
is on it will try to wield a launcher matching the ammo in the quiver.

%.pg
Some characters have the ability to throw or shoot a volley of multiple
items (from the same stack) in a single action.
Knowing how to load several rounds of ammunition at
once---or hold several missiles in your hand---and
still hit a target is not an easy task.
Rangers are among those who are adept
at this task, as are those with a high level of proficiency in the
relevant weapon skill (in bow skill if you're wielding one to
shoot arrows, in crossbow skill if you're wielding one to shoot bolts,
or in sling skill if you're wielding one to shoot stones).
The number of items that the character has a chance to fire varies from
turn to turn.  You can explicitly limit the number of shots by using a
numeric prefix before the `\texttt{t}' or `\texttt{f}' command.
For example, ``\texttt{2f}'' (or ``\texttt{n2f}'' if using
\textit{number\textunderscore pad}
mode) would ensure that at most 2 arrows are shot
even if you could have fired 3.  If you specify
a larger number than would have been shot (``\texttt{4f}'' in this example),
you'll just end up shooting the same number (3, here) as if no limit
had been specified.  Once the volley is in motion, all of the items
will travel in the same direction; if the first ones kill a monster,
the others can still continue beyond that spot.

%.hn 3
\subsection*{Weapon proficiency}

%.pg
You will have varying degrees of skill in the weapons available.
Weapon proficiency, or weapon skills, affect how well you can use
particular types of weapons, and you'll be able to improve your skills
as you progress through a game, depending on your role, your experience
level, and use of the weapons.

%.pg
For the purposes of proficiency, weapons have
been divided up into various groups such as daggers, broadswords, and
polearms.  Each role has a limit on what level of proficiency a character
can achieve for each group.  For instance, wizards can become highly
skilled in daggers or staves but not in swords or bows.

%.pg
The ``\texttt{\#enhance}'' extended command is used to review current weapons
proficiency
(also spell proficiency) and to choose which skill(s) to improve when
you've used one or more skills enough to become eligible to do so.  The
skill rankings are ``none'' (sometimes also referred to as ``restricted'',
because you won't be able to advance), ``unskilled'', ``basic'', ``skilled'',
and ``expert''.  Restricted skills simply will not appear in the list
shown by ``\texttt{\#enhance}''.
(Divine intervention might unrestrict a particular
skill, in which case it will start at unskilled and be limited to basic.)
Some characters can enhance their barehanded combat or martial arts skill
beyond expert to ``master'' or ``grand master''.

%.pg
Use of a weapon in which you're restricted or unskilled
will incur a modest penalty in the chance to hit a monster and also in
the amount of damage done when you do hit; at basic level, there is no
penalty or bonus; at skilled level, you receive a modest bonus in the
chance to hit and amount of damage done; at expert level, the bonus is
higher.  A successful hit has a chance to boost your training towards
the next skill level (unless you've already reached the limit for this
skill).  Once such training reaches the threshold for that next level,
you'll be told that you feel more confident in your skills.  At that
point you can use ``\texttt{\#enhance}'' to increase one or more skills.
Such skills
are not increased automatically because there is a limit to your total
overall skills, so you need to actively choose which skills to enhance
and which to ignore.

%.hn 3
\subsection*{Two-Weapon combat}

%.pg
Some characters can use two weapons at once.  Setting things up to
do so can seem cumbersome but becomes second nature with use.
To wield two weapons, you need to use the ``\texttt{\#twoweapon}'' command.
But first you need to have a weapon in each hand.
(Note that your two weapons are not fully equal; the one in the
hand you normally wield with is considered primary and the other
one is considered secondary.
The most noticeable difference is after you
stop---or before you begin, for that matter---wielding
two weapons at once.
The primary is your wielded weapon and the
secondary is just an item in your inventory that's been designated
as alternate weapon.)

%.pg
If your primary weapon is wielded but your off hand is empty or has
the wrong weapon, use the sequence `\texttt{x}', `\texttt{w}', `\texttt{x}' to
first swap your
primary into your off hand, wield whatever you want as secondary
weapon, then swap them both back into the intended hands.
If your secondary or alternate weapon is correct but your primary
one is not, simply use `\texttt{w}' to wield the primary.
Lastly, if neither hand holds the correct weapon,
use `\texttt{w}', `\texttt{x}', `\texttt{w}'
to first wield the intended secondary, swap it to off hand, and then
wield the primary.

%.pg
The whole process can be simplified via use of the
\textit{pushweapon}
option.  When it is enabled, then using `\texttt{w}' to wield something
causes the currently wielded weapon to become your alternate weapon.
So the sequence `\texttt{w}', `\texttt{w}' can be used to first wield the weapon you
intend to be secondary, and then wield the one you want as primary
which will push the first into secondary position.

%.pg
When in two-weapon combat mode, using the `\texttt{X}' command
toggles back to single-weapon mode.
Throwing or dropping either of the
weapons or having one of them be stolen or destroyed will also make you
revert to single-weapon combat.

%.hn 2
\subsection*{Armor (`\texttt{[}')}

%.pg
Lots of unfriendly things lurk about; you need armor to protect
yourself from their blows.  Some types of armor offer better
protection than others.  Your armor class is a measure of this
protection.  Armor class (AC) is measured as in AD\&D, with 10 being
the equivalent of no armor, and lower numbers meaning better armor.
Each suit of armor which exists in AD\&D gives the same protection in
\textit{NetHack}.

Here is a list of the armor class values provided by suits of armor:

\begin{center}
\begin{tabular}{lllll}
dragon scale mail      & 1 & \makebox[20mm]{}  & plate mail            & 3\\
crystal plate mail     & 3 &                   & bronze plate mail     & 4\\
splint mail            & 4 &                   & banded mail           & 4\\
dwarvish mithril-coat  & 4 &                   & elven mithril-coat    & 5\\
chain mail             & 5 &                   & orcish chain mail     & 6\\
scale mail             & 6 &                   & dragon scales         & 7\\
studded leather armor  & 7 &                   & ring mail             & 7\\
orcish ring mail       & 8 &                   & leather armor         & 8\\
leather jacket         & 9 &                   & no armor              & 10\\
\end{tabular}
\end{center}

%.pg
\noindent You can also wear other pieces of armor (cloak over suit, shirt under
suit, helmet, gloves, boots, shield) to lower your armor class even
further.
%--too obvious to mention unless we include polymorph into ettin or maralith
% You can wear at most one item of each category (one suit of armor, one
% cloak, one helmet, one shield, and so on) at a time.
Most of these provide a one or two point improvement to AC (making the
overall value smaller and eventually negative) but can also be
enchanted.
Shirts are an exception; they don't provide any protection unless enchanted.
Some cloaks also don't improve AC when unenchanted but all cloaks offer
some protection against rust or corrosion to suits worn under them and
against some monster \textit{touch} attacks.

%.pg
If a piece of armor is enchanted, its armor protection will be better
(or worse) than normal, and its ``plus'' (or minus) will subtract from
your armor class.  For example, a +1 chain mail would give you
better protection than normal chain mail, lowering your armor class one
unit further to 4.  When you put on a piece of armor, you immediately
find out the armor class and any ``plusses'' it provides.  Cursed
pieces of armor usually have negative enchantments (minuses) in
addition to being unremovable.

%.pg
Many types of armor are subject to some kind of damage like rust.  Such
damage can be repaired.  Some types of armor may inhibit spell casting.

%.pg
The \textit{nudist}
option can be set (prior to game start) to attempt to play the entire
game without wearing any armor (a self-imposed challenge which is
extremely difficult to accomplish).

%.pg
Some interfaces support the \texttt{armorstatus} option.
When it is enabled, an extra status condition is displayed, summarizing
currently worn armor.

%.pg
The commands to use armor are `\texttt{W}' (wear) and `\texttt{T}' (take off).
The `\texttt{A}' command can be used to take off armor as well as other
worn items.
Also, `\texttt{P}' (put on) and `\texttt{R}' (remove) which are normally for
accessories can be used for armor, but pieces of armor won't be shown
as likely candidates in a prompt for choosing what to put on or remove.

%.hn 2
\subsection*{Food (`\texttt{\%}')}

%.pg
Food is necessary to survive.  If you go too long without eating you
will faint, and eventually die of starvation.
Some types of food will spoil, and become unhealthy to eat,
if not protected.
Food stored in ice boxes or tins (``cans'')
will usually stay fresh, but ice boxes are heavy, and tins
take a while to open.

%.pg
When you kill monsters, they usually leave corpses which are also
``food.''  Many, but not all, of these are edible; some also give you
special powers when you eat them.  A good rule of thumb is ``you are
what you eat.''

%.pg
Some character roles and some monsters are vegetarian.  Vegetarian monsters
will typically never eat animal corpses, while vegetarian players can,
but with some rather unpleasant side-effects.

%.pg
You can name one food item after something you like to eat with the
\textit{fruit} option.

%.pg
The command to eat food is `\texttt{e}'.

%.hn 2
\subsection*{Scrolls (`\texttt{?}')}

%.pg
Scrolls are labeled with various titles, probably chosen by ancient wizards
for their amusement value (for example, ``READ ME,'' or ``THANX MAUD'' backwards).
Scrolls disappear after you read them (except for blank ones, without
magic spells on them).

%.pg
One of the most useful of these is the %
\textit{scroll of identify}, which
can be used to determine what another object is, whether it is cursed or
blessed, and how many uses it has left.
Some objects of subtle enchantment are difficult to identify without these.

%.pg
A scroll whose label is known can be read even when the hero is blind.
If a scroll has been discovered, it will be listed in inventory by type
rather than by label, but the label is known in that situation even though
it isn't shown.

%.pg
Many scrolls produce a different effect from usual if they are blessed or
cursed, or read while the hero is confused.

%.pg
A mail daemon may run up and deliver mail to you as a %
\textit{scroll of mail} (on versions compiled with this feature).
To use this feature on versions where \textit{NetHack}
mail delivery is triggered by electronic mail appearing in your system mailbox,
you must let \textit{NetHack} know where to look for new mail by setting the
\texttt{MAIL} environment variable to the file name of your mailbox.
You may also want to set the \texttt{MAILREADER} environment variable to the
file name of your favorite reader, so \textit{NetHack} can shell to it when you
read the scroll.
On versions of \textit{NetHack} where mail is randomly
generated internal to the game, these environment variables are ignored.
You can disable the mail daemon by turning off the
\textit{mail} option.

%.pg
The command to read a scroll is `\texttt{r}'.

%.hn 2
\subsection*{Potions (`\texttt{!}')}

%.pg
Potions are distinguished by the color of the liquid inside the flask.
They disappear after you quaff them.

%.pg
Clear potions are potions of water.  Sometimes these are
blessed or cursed, resulting in holy or unholy water.  Holy water is
the bane of the undead, so potions of holy water are good things to
throw (`\texttt{t}') at them.  It is also sometimes very useful to dip
(``\texttt{\#dip}'') an object into a potion.

%.pg
The command to drink a potion is `\texttt{q}' (quaff).

%.hn 2
\subsection*{Wands (`\texttt{/}')}

%.pg
Wands usually have multiple magical charges.
Some types of wands require a direction in which to zap them.
You can also
zap them at yourself (just give a `\texttt{.}' or `\texttt{s}' for the direction).
Be warned, however, for this is often unwise.
Other types of wands
don't require a direction.  The number of charges in a
wand is random and decreases by one whenever you use it.

%.pg
When the number of charges left in a wand becomes zero, attempts to use the
wand will usually result in nothing happening.  Occasionally, however, it may
be possible to squeeze the last few mana points from an otherwise spent wand,
destroying it in the process.  A wand may be recharged by using suitable
magic, but doing so runs the risk of causing it to explode.  The chance
for such an explosion starts out very small and increases each time the
wand is recharged.

%.pg
In a truly desperate situation, when your back is up against the wall, you
might decide to go for broke and break your wand.  This is not for the faint
of heart.  Doing so will almost certainly cause a catastrophic release of
magical energies.

%.pg
When you have fully identified a particular wand, inventory display will
include additional information in parentheses: the number of times it has
been recharged followed by a colon and then by its current number of charges.
A current charge count of \texttt{-1} is a special case indicating that the wand
has been cancelled.

%.pg
The command to use a wand is `\texttt{z}' (zap).  To break one, use the `\texttt{a}'
(apply) command.

%.hn 2
\subsection*{Rings (`\texttt{=}')}

%.pg
Rings are very useful items, since they are relatively permanent
magic, unlike the usually fleeting effects of potions, scrolls, and
wands.

%.pg
Putting on a ring activates its magic.
You can wear at most two rings at any time, one on the ring finger of
each hand.

%.pg
Most worn rings also cause you to grow hungry more rapidly, the rate
varying with the type of ring.

%.pg
When wearing gloves, rings are worn underneath.
If the gloves are cursed, rings cannot be put on and any already being
worn cannot be removed.
When worn gloves aren't cursed, you don't have to manually take them
off before putting on or removing a ring and then re-wear them after.
That's done implicitly to avoid unnecessary tedium.

%.pg
The commands to use rings are `\texttt{P}' (put on) and `\texttt{R}' (remove).
`\texttt{A}', `\texttt{W}', and `\texttt{T}' can also be used; see \textit{Amulets}.

%.hn 2
\subsection*{Spellbooks (`\texttt{+}')}

%.pg
Spellbooks are tomes of mighty magic.  When studied with the `\texttt{r}' (read)
command, they transfer to the reader the knowledge of a spell (and
therefore eventually become
unreadable)---unless the attempt backfires.
Reading a cursed spellbook or one with mystic runes beyond
your ken can be harmful to your health!

%.pg
A spell (even when learned) can also backfire when you cast it.  If you
attempt to cast a spell well above your experience level, or if you have
little skill with the appropriate spell type, or cast it at
a time when your luck is particularly bad, you can end up wasting both the
energy and the time required in casting.

%.pg
Casting a spell calls forth magical energies and focuses them with
your naked mind.  Some of the magical energy released comes from within
you.
Casting temporarily drains your magical power, which will slowly be
recovered, and causes you to need additional food.
Casting of spells also requires practice.  With practice, your
skill in each category of spell casting will improve.  Over time, however,
your memory of each spell will dim, and you will need to relearn it.

%.pg
Some spells require a direction in which to cast them, similar to wands.
To cast one at yourself, just give a `\texttt{.}' or `\texttt{s}' for the direction.
A few spells require you to pick a target location rather than just specify
a particular direction.
Other spells don't require any direction or target.

%.pg
Just as weapons are divided into groups in which a character can become
proficient (to varying degrees), spells are similarly grouped.
Successfully casting a spell exercises its skill group; using the
``\texttt{\#enhance}'' command to advance a sufficiently exercised skill
will affect all spells within the group.  Advanced skill may increase the
potency of spells, reduce their risk of failure during casting attempts,
and improve the accuracy of the estimate for how much longer they will
be retained in your memory.
Skill slots are shared with weapons skills.  (See also the section on
``Weapon proficiency''.)

%.pg
Casting a spell also requires flexible movement, and wearing various types
of armor may interfere with that.

%.pg
The command to read a spellbook is the same as for scrolls, `\texttt{r}' (read).
The `\texttt{+}' command lists each spell you know along with its level, skill
category, chance of failure when casting, and an estimate of how strongly
it is remembered.
The `\texttt{Z}' (cast) command casts a spell.

%.hn 2
\subsection*{Tools (`\texttt{(}')}

%.pg
Tools are miscellaneous objects with various purposes.  Some tools
have a limited number of uses, akin to wand charges.  For example, lamps burn
out after a while.  Other tools are containers, which objects can
be placed into or taken out of.

%.pg
Some tools (such as a blindfold) can be \textit{worn} and can be put on and
removed like other accessories (rings, amulets); see \textit{Amulets}.
Other tools (such as pick-axe) can be wielded as weapons in addition to
being applied for their usual purpose, and in some cases (again, pick-axe)
become wielded as a weapon even when applied.

%.pg
% Mentioned here because of the old method of attempting "Zen" conduct:
% restart until there's a blindfold in starting inventory and put it on
% first thing.
The \textit{blind}
option can be set (prior to game start) to attempt to play the entire
game without being able to see (a self-imposed challenge which is
very difficult to accomplish).

%.pg
The command to use a tool is `\texttt{a}' (apply).

%.hn 3
\subsection*{Containers}

%.pg
You may encounter bags, boxes, and chests in your travels.  A tool of
this sort can be opened with the ``\texttt{\#loot}'' extended command when
you are standing on top of it (that is, on the same floor spot),
or with the `\texttt{a}' (apply) command when you are carrying it.  However,
chests are often locked, and are in any case unwieldy objects.
You must set one down before unlocking it by
using a key or lock-picking tool with the `\texttt{a}' (apply) command,
by kicking it with the `\texttt{\textasciicircum D}' command,
or by using a weapon to force the lock with the ``\texttt{\#force}''
extended command.

%.pg
Some chests are trapped, causing nasty things to happen when you
unlock or open them.  You can check for and try to deactivate traps
with the ``\texttt{\#untrap}'' extended command.

%.pg
When the contents of a container are known, that container will be
described as something like ``a sack containing 3 items''.
In this example, the 3 refers to number of \textit{stacks} of compatible
items, not to the total number of individual items.
So a sack holding 2 sky blue potions, 7 arrows, and 350 gold pieces would be
described as having 3 items rather than 10 or 359.
And you would need to have 3 unused inventory slots available in order
to take everything out (for the case where the items you remove don't
combine into bigger stacks with things you're already carrying).

%.pg
If a chest or large box is described as ``broken'', that means that it
can't be locked rather than that it no longer functions as a container.

%.pg
The \textit{apply} and \textit{loot} commands allow you to take out and/or
put in an arbitrary number of items in a single operation.
If you want to take everything out of a container, you can use the
``\texttt{\#tip}'' command to pour the contents onto the floor.
This may be your only way to get things out if your hands are stuck
to a cursed two-handed weapon.
When your hands aren't stuck, you have the potential to pour the
contents into another container.
(As of this writing, the other container must be carried rather than on
the floor.)

%.hn 2
\subsection*{Amulets (`\texttt{"}')}

%.pg
Amulets are very similar to rings, and often more powerful.  Like
rings, amulets have various magical properties, some beneficial,
some harmful, which are activated by putting them on.

%.pg
Only one amulet may be worn at a time, around your neck.
Like wearing rings, wearing an amulet affects your metabolism, causing
you to grow hungry more rapidly.

%.pg
The commands to use amulets are the same as for rings, `\texttt{P}' (put on)
and `\texttt{R}' (remove).
`\texttt{A}' can be used to remove various worn items including amulets.
Also, '\texttt{W}' (wear) and `\texttt{T}' (take off) which are normally for
armor can be used for amulets and other accessories (rings and eyewear),
but accessories won't be shown as likely candidates in a prompt for
choosing what to wear or take off.

%.hn 2
\subsection*{Gems (`\texttt{*}')}

%.pg
Some gems are valuable, and can be sold for a lot of gold.
They are also a far more efficient way of carrying your riches.
Valuable gems increase your score if you bring them with you when you exit.

%.pg
Other small rocks are also categorized as gems, but they are much less
valuable.
All rocks, however, can be used as projectile weapons (if you have a sling).
In the most desperate of cases, you can still throw them by hand.

%.hn 2
\subsection*{Large rocks (`\texttt{`}')}
%.pg
Statues and boulders are not particularly useful, and are generally heavy.
It is rumored that some statues are not what they seem.

%.pg
Boulders occasionally block your path.
You can push one forward (by attempting to walk onto its spot)
when nothing blocks \textit{its} path, or you can
smash it into a pile of small rocks with breaking magic or a pick-axe.
It is possible to move onto a boulder's location if certain conditions
are met; ordinarily one of those conditions is that pushing it any
further be blocked.
Using the move-without-picking-up prefix (default key `\texttt{m}')
prior to the direction of movement will attempt to move to a boulder's
location without pushing it in addition to the prefix's usual action of
suppressing auto-pickup at the destination.

%.pg
Very large humanoids (giants and their ilk) have been known to pick up
boulders and use them as missile weapons.

%.pg
Unlike boulders, statues can't be pushed, but don't need to be because
they don't block movement.
%.\" 'rumor' above is about statue traps; this is a hint about statue contents
They can be smashed into rocks though.

%.pg
For some configurations of the program, statues are no longer shown
as `\texttt{`}'
but by the letter representing the monster they depict instead.

%.hn 2
\subsection*{Gold (`\texttt{\$}')}

%.pg
Gold adds to your score, and you can buy things in shops with it.
There are a number
of monsters in the dungeon that may be influenced by the amount of gold
you are carrying (shopkeepers aside).

%.pg
Gold pieces are the only type of object where bless/curse state does not
apply.
They're always uncursed but never described as uncursed even if you turn
off the ``\textit{implicit\textunderscore uncursed}'' option.
You can set the ``\textit{goldX}''
option if you prefer to have gold pieces be treated as bless/curse state
\textit{unknown} rather than as known to be uncursed.
Only matters when you're using an object selection prompt that can filter
by ``\texttt{BUCX}'' state.

%.hn 2
\subsection*{Persistence of Objects}

%.pg
Normally, if you have seen an object at a particular map location and
move to another location where you can't directly see that object any
more, it will continue to be displayed on your map.
That remains the case even if it is not actually there any
more---perhaps a monster has picked it up or it has rotted away---until
you can see or feel that location again.
One notable exception is that if the object gets covered by the
``remembered, unseen monster'' marker.
When that marker is later removed
after you've verified that no monster is there, you will have forgotten that
there was any object there regardless of whether the unseen monster
actually took the object.
If the object is still there, then once you see or feel that location
again you will re-discover the object and resume remembering it.

%.pg
The situation is the same for a pile of objects, except that only the
top item of the pile is displayed.
The
\textit{hilite\textunderscore pile}
option can be enabled in order to show an item differently when it is
the top one of a pile.

%.hn 1
\section{Conduct}

%.pg
As if winning \textit{NetHack} were not difficult enough, certain players
seek to challenge themselves by imposing restrictions on the
way they play the game.  The game automatically tracks some of
these challenges, which can be checked at any time with the \texttt{\#conduct}
command or at the end of the game.  When you perform an action which
breaks a challenge, it will no longer be listed.  This gives
players extra ``bragging rights'' for winning the game with these
challenges.  Note that it is perfectly acceptable to win the game
without resorting to these restrictions and that it is unusual for
players to adhere to challenges the first time they win the game.

%.pg
Several of the challenges are related to eating behavior.  The most
difficult of these is the foodless challenge.  Although creatures
can survive long periods of time without food, there is a physiological
need for water; thus there is no restriction on drinking beverages,
even if they provide some minor food benefits.
Calling upon your god for help with starvation does
not violate any food challenges either.

%.pg
A strict vegan diet is one which avoids any food derived from animals.
The primary source of nutrition is fruits and vegetables.  The
corpses and tins of blobs (`\texttt{b}'), jellies (`\texttt{j}'), and fungi
(`\texttt{F}') are also considered to be vegetable matter.  Certain human
food is prepared without animal products; namely, lembas wafers, cram
rations, food rations (gunyoki), K-rations, and C-rations.
Metal or another normally indigestible material eaten while polymorphed
into a creature that can digest it is also considered vegan food.
Note however that eating such items still counts against foodless conduct.

%.pg
Vegetarians do not eat animals;
however, they are less selective about eating animal byproducts than vegans.
In addition to the vegan items listed above, they may eat any kind
of pudding (`\texttt{P}') other than the black puddings,
eggs and food made from eggs (fortune cookies and pancakes),
food made with milk (cream pies and candy bars), and lumps of
royal jelly.  Monks are expected to observe a vegetarian diet.

%.pg
Eating any kind of meat violates the vegetarian, vegan, and foodless
conducts.  This includes tripe rations, the corpses or tins of any
monsters not mentioned above, and the various other chunks of meat
found in the dungeon.  Swallowing and digesting a monster while polymorphed
is treated as if you ate the creature's corpse.
Eating leather, dragon hide, or bone items while
polymorphed into a creature that can digest it, or eating monster brains
while polymorphed into a mind flayer, is considered eating
an animal, although wax is only an animal byproduct.

%.pg
Regardless of conduct, there will be some items which are indigestible,
and others which are hazardous to eat.  Using a swallow-and-digest
attack against a monster is equivalent to eating the monster's corpse.
Please note that the term ``vegan'' is used here only in the context of
diet.  You are still free to choose not to use or wear items derived
from animals (e.g. leather, dragon hide, bone, horns, coral), but the
game will not keep track of this for you.  Also note that ``milky''
potions may be a translucent white, but they do not contain milk,
so they are compatible with a vegan diet.  Slime molds or
player-defined ``fruits'', although they could be anything
from ``cherries'' to ``pork chops'', are also assumed to be vegan.

%.pg
An atheist is one who rejects religion.  This means that you cannot
\texttt{\#pray}, \texttt{\#offer} sacrifices to any god,
\texttt{\#turn} undead, or \texttt{\#chat} with a priest.
Particularly selective readers may argue that playing Monk or Priest
characters should violate this conduct; that is a choice left to the
player.  Offering the Amulet of Yendor to your god is necessary to
win the game and is not counted against this conduct.  You are also
not penalized for being spoken to by an angry god, priest(ess), or
other religious figure; a true atheist would hear the words but
attach no special meaning to them.

%.pg
A pauper starts the game with no possessions, no spells, and no weapon or
spell skills (and if playing as a knight, your pony will not have a saddle).
Can only be initiated by starting a new game with \texttt{OPTIONS=pauper}
set in your run-time configurtion file or \texttt{NETHACKOPTIONS} environment
variable.
Once the game is underway, you can acquire and use items, spells, and skills
in the usual way.

%.pg
Most players fight with a wielded weapon (or tool intended to be
wielded as a weapon).  Another challenge is to win the game without
using such a wielded weapon.  You are still permitted to throw,
fire, and kick weapons; use a wand, spell, or other type of item;
or fight with your hands and feet.

%.pg
In \textit{NetHack}, a pacifist refuses to cause the death of any other monster
(i.e. if you would get experience for the death).  This is a particularly
difficult challenge, although it is still possible to gain experience
by other means.

%.pg
An illiterate character does not read or write.  This includes reading
a scroll, spellbook, fortune cookie message, or t-shirt; writing a
scroll; or making an engraving of anything other than a single ``X'' (the
traditional signature of an illiterate person).  Reading an engraving,
or any item that is absolutely necessary to win the game, is not counted
against this conduct.  The identity of scrolls and spellbooks (and
knowledge of spells) in your starting inventory is assumed to be
learned from your teachers prior to the start of the game and isn't
counted.

%.pg
There is a side-branch to the main dungeon called ``Sokoban,'' briefly
described in the earlier section about \textit{Traps}.
As mentioned there, the goal is to push boulders into pits and/or holes
to plug those in order to both get the boulders out of your way and be
able to go past the traps.
There are some special ``rules'' that are active when in that branch
of the dungeon.
Some rules can't be bypassed, such as being unable to push a boulder
diagonally.
Other rules can, such as not smashing boulders with magic or tools,
but doing so causes you to receive a luck penalty.
No message about that is given at the time, but it is tracked as a conduct.
The \texttt{\#conduct} command and end of game disclosure will report whether
you have abided by the special rules of Sokoban, and if not, how many
times you violated them, providing you with a way to discover which
actions incur bad luck so that you can be better informed about whether
or not to avoid repeating those actions in the future.
(Note:  the Sokoban conduct will only be displayed if you have
entered the Sokoban branch of the dungeon during the current game.
Once that has happened, it becomes part of disclosed conduct even if
you haven't done anything interesting there.
Ending the game with ``never broke the Sokoban rules'' conduct is most
meaningful if you also manage to perform
the ``obtained the Sokoban prize'' achievement
(see \textit{Achievements} below).)

%.pg
There are several other challenges tracked by the game.  It is possible
to eliminate one or more species of monsters by genocide; playing without
this feature is considered a challenge.  When the game offers you an
opportunity to genocide monsters, you may respond with the monster type
``none'' if you want to decline.  You can change the form of an item into
another item of the same type (``polypiling'') or the form of your own
body into another creature (``polyself'') by wand, spell, or potion of
polymorph; avoiding these effects are each considered challenges.
Polymorphing monsters, including pets, does not break either of these
challenges.
Finally, you may sometimes receive wishes; a game without an attempt to
wish for any items is a challenge, as is a game without wishing for
an artifact (even if the artifact immediately disappears).  When the
game offers you an opportunity to make a wish for an item, you may
choose ``nothing'' if you want to decline.

%.hn 2
\subsection*{Achievements}

%.pg
End of game disclosure will also display various achievements
representing progress toward ultimate ascension, if any have been
attained.
They aren't directly related to \textit{conduct} but are grouped with
it because they fall into the same category of ``bragging rights''
and to limit the number of questions during disclosure.
Listed here roughly in order of difficulty and not necessarily in the order
in which you might accomplish them.

% [length stuff copied from paranoid_confirmation]
\newlength{\achwidth}
%.PS "Mines'\~End\~"
\settowidth{\achwidth}{\texttt{Mines'~End~}}
\addtolength{\achwidth}{\labelsep}
\begin{description}[leftmargin=\achwidth, topsep=1mm, itemsep=0mm, labelwidth=*, font=\ttfamily, align=right]
%.PL Shop
\item[<Rank>]
Attained rank title \textit{Rank}.
\item[Shop]
Entered a shop.
\item[Temple]
Entered a temple.
\item[Mines]
Entered the Gnomish Mines.
\item[Town]
Entered Mine Town.
\item[Oracle]
Consulted the Oracle of Delphi.
\item[Novel]
Read a passage from a Discworld Novel.
\item[Sokoban]
Entered Sokoban.
\item["Big~Room"]
Entered the Big Room.
\item["Soko-Prize"]
Explored to the top of Sokoban and found a special item there.
\item[Mines'~End]
Explored to the bottom of the Gnomish Mines and found a special item there.
\item[Medusa]
Defeated Medusa.
\item[Tune]
Discovered the tune that can be used to open and close the drawbridge on
the Castle level.
\item[Bell]
Acquired the Bell of Opening.
\item[Gehennom]
Entered Gehennom.
\item[Candle]
Acquired the Candelabrum of Invocation.
\item[Book]
Acquired the Book of the Dead.
\item[Invocation]
Gained access to the bottommost level of Gehennom.
\item[Amulet]
Acquired the fabled Amulet of Yendor.
\item[Endgame]
Reached the Elemental Planes.
\item[Astral]
Reached the Astral Plane level.
\item[Blind]
Blind from birth.
\item[Deaf]
Deaf from birth.
\item[Nudist]
Never wore any armor.
\item[Pauper]
Started out with no possessions.
\item[Ascended]
Delivered the Amulet to its final destination.
\end{description}
%.PE
%.ED

%.lp "Notes:  "
\noindent
Notes:

%.pg
Achievements are recorded and subsequently reported in the order in which
they happen during your current game rather than the order listed here.

%.pg
There are nine \textit{Rank} titles for each role, bestowed at experience
levels 1, 3, 6, 10, 14, 18, 22, 26, and 30.
The one for experience level 1 is not recorded as an achievement.
Losing enough levels to revert to lower rank(s) does not discard the
corresponding achievement(s).

%.pg
There's no guaranteed \textit{Novel} so the achievement to read one might
not always be attainable (except perhaps by wishing).
Similarly, the \textit{Big Room} level is not always present.
Unlike with the Novel, there's no way to wish for this opportunity.

%.pg
The ``special items'' hidden in \textit{Mines'~End} and \textit{Sokoban}
are not unique but are considered to be prizes or rewards
for exploring those levels since doing so is not necessary to complete
the game.
Finding other instances of the same objects doesn't record the
corresponding achievement.

%.pg
The \textit{Medusa} achievement is recorded if she dies for any reason,
even if you are not directly responsible, and only if she dies.

%.pg
The 5-note \textit{tune} can be learned via trial and error with a musical
instrument played closely
enough---but not too close!---to
the Castle level's drawbridge or can be given to you via prayer boon.

%.pg
\textit{Blind}, \textit{Deaf}, \textit{Nudist}, and \textit{Pauper} are also conducts, and they can only be
enabled by setting the correspondingly named option in \texttt{NETHACKOPTIONS}
or run-time configuration file prior to game start.
In the case of \textit{Blind} and \textit{Deaf}, the option also enforces the
conduct.
They aren't really significant accomplishments unless/until you make
substantial progress into the dungeon.

%.hn 1
%.hn 1
\section{Scoring}

%.pg
\textit{NetHack} maintains a list of the top scores or scorers on your machine,
depending on how it is set up.  In the latter case, each account on
the machine can post only one non-winning score on this list.  If
you score higher than someone else on this list, or better your
previous score, you will be inserted in the proper place under your
current name.  How many scores are kept can also be set up when
\textit{NetHack} is compiled.

%.pg
Your score is chiefly based upon how much experience you gained, how
much loot you accumulated, how deep you explored, and how the game
ended.  If you quit the game, you escape with all of your gold intact.
If, however, you get killed in the Mazes of Menace, the guild will
only hear about 90\,\% of your gold when your corpse is discovered
(adventurers have been known to collect finder's fees).  So, consider
whether you want to take one last hit at that monster and possibly
live, or quit and stop with whatever you have.  If you quit, you keep
all your gold, but if you swing and live, you might find more.

%.pg
If you just want to see what the current top players/games list is, you
can type
\begin{verbatim}
    nethack -s all
\end{verbatim}
on most versions.

%.hn 1
\section{Explore mode}

%.pg
\textit{NetHack} is an intricate and difficult game.  Novices might falter
in fear, aware of their ignorance of the means to survive.  Well, fear
not.  Your dungeon comes equipped with an ``explore'' or ``discovery''
mode that enables you to keep old save files and cheat death, at the
paltry cost of not getting on the high score list.

%.pg
There are two ways of enabling explore mode.  One is to start the game
with the \texttt{-X}
command-line switch or with the
\textit{playmode:explore}
option.  The other is to issue the `\texttt{\#exploremode}' extended command while
already playing the game.  Starting a new game in explore mode provides your
character with a wand of wishing in initial inventory; switching
during play does not.  The other benefits of explore mode are left for
the trepid reader to discover.

%.pg
%.hn 2
\subsection*{Debug mode}

%.pg
Debug mode, also known as wizard mode, is undocumented aside from this
brief description and the various ``debug mode only'' commands listed
among the command descriptions.
It is intended for tracking down problems within the
program rather than to provide god-like powers to your character, and
players who attempt debugging are expected to figure out how to use it
themselves.
It is initiated by starting the game with the
\texttt{-D}
command-line switch or with the
\textit{playmode:debug}
option.

%.pg
For some systems, the player must be logged in
under a particular user name to be allowed to use debug mode; for others,
the hero must be given a particular character name (but may be any role;
there's no connection between ``wizard mode'' and the \textit{Wizard} role).
Attempting to start a game in debug mode when not allowed
or not available will result in falling back to explore mode instead.

%.hn
\section{Credits}
%.pg
The original %
\textit{hack} game was modeled on the Berkeley
%.ux
UNIX
\textit{rogue} game.  Large portions of this document were shamelessly
cribbed from %
\textit{A Guide to the Dungeons of Doom}, by Michael C. Toy
and Kenneth C. R. C. Arnold.  Small portions were adapted from
\textit{Further Exploration of the Dungeons of Doom}, by Ken Arromdee.

%.pg
\textit{NetHack} is the product of literally scores of people's work.
Main events in the course of the game development are described below:

%.pg
\bigskip
\noindent \textit{Jay Fenlason} wrote the original \textit{Hack}, with help from {\it
Kenny Woodland}, \textit{Mike Thome}, and \textit{Jon Payne}.

%.pg
\medskip
\noindent \textit{Andries Brouwer} did a major re-write while at
Stichting Mathematisch Centrum (now Centrum Wiskunde \& Informatica),
transforming Hack into a very different game.
He published the Hack source code for use on UNIX
systems by posting that to Usenet
newsgroup \textit{net.sources} (later renamed \textit{comp.sources})
releasing version 1.0 in December of 1984, then versions 1.0.1, 1.0.2,
and finally 1.0.3 in July of 1985.
Usenet newsgroup \textit{net.games.hack} (later
renamed \textit{rec.games.hack}, eventually replaced
by \textit{rec.games.roguelike.nethack})
was created for discussing it.

%.pg
\medskip
\noindent \textit{Don G. Kneller} ported \textit{Hack} 1.0.3 to Microsoft C and MS-DOS,
producing \textit{PC Hack} 1.01e, added support for DEC Rainbow graphics in
version 1.03g, and went on to produce at least four more versions (3.0, 3.2,
3.51, and 3.6;
note that these are old \textit{Hack} version numbers, not contemporary
\textit{NetHack} ones).

%.pg
\medskip
\noindent \textit{R. Black} ported \textit{PC Hack} 3.51 to Lattice C and the Atari
520/1040ST, producing \textit{ST Hack} 1.03.

%.pg
\medskip
\noindent \textit{Mike Stephenson} merged these various versions back together,
incorporating many of the added features, and produced \textit{NetHack} version
1.4 in 1987.
He then coordinated a cast of thousands in enhancing and debugging
\textit{NetHack} 1.4 and released \textit{NetHack} versions 2.2 and 2.3.
Like Hack, they were released by posting their source code to Usenet where
they remained available in various archives accessible
via \textit{ftp} and \textit{uucp} after expiring from the newsgroup.

%.pg
\medskip
\noindent Later, Mike coordinated a major re-write of the game, heading a team which
included \textit{Ken Arromdee}, \textit{Jean-Christophe Collet}, \textit{Steve Creps},
\textit{Eric Hendrickson}, \textit{Izchak Miller}, \textit{Eric S. Raymond}, \textit{John
Rupley}, \textit{Mike Threepoint}, and \textit{Janet Walz}, to produce
\textit{NetHack} 3.0c.

%.pg
\medskip
\noindent \textit{NetHack} 3.0 was ported to the Atari by \textit{Eric R. Smith}, to OS/2 by
\textit{Timo Hakulinen}, and to VMS by \textit{David Gentzel}.  The three of them
and \textit{Kevin Darcy} later joined the main \textit{NetHack Development Team} to produce
subsequent revisions of 3.0.

%.pg
\medskip
\noindent \textit{Olaf Seibert} ported \textit{NetHack} 2.3 and 3.0 to the Amiga.  {\it
Norm Meluch}, \textit{Stephen Spackman} and \textit{Pierre Martineau} designed
overlay code for \textit{PC NetHack} 3.0.  \textit{Johnny Lee} ported
\textit{NetHack} 3.0 to the Macintosh.  Along with various other Dungeoneers, they
continued to enhance the PC, Macintosh, and Amiga ports through the later
revisions of 3.0.

%.pg
Version 3.0 went through ten relatively rapidly released ``patch-level''
revisions.
Versions at the time were known as 3.0 for the base release and variously
as ``3.0a'' through ``3.0j'',
``3.0~patchlevel~1'' through ``3.0~patchlevel~10'',
or ``3.0pl1'' through ``3.0pl10''
rather than 3.0.0 and 3.0.1 through 3.0.10;
the three component numbering scheme began to be used with 3.1.0.

%.pg
\medskip
\noindent Headed by \textit{Mike Stephenson} and coordinated by \textit{Izchak Miller}
and \textit{Janet Walz}, the \textit{NetHack Development Team} which now included
\textit{Ken Arromdee},
\textit{David Cohrs}, \textit{Jean-Christophe Collet}, \textit{Kevin Darcy},
\textit{Matt Day}, \textit{Timo Hakulinen}, \textit{Steve Linhart}, \textit{Dean Luick},
\textit{Pat Rankin}, \textit{Eric Raymond}, and \textit{Eric Smith} undertook a
radical revision of 3.0.
They re-structured the game's design, and re-wrote major
parts of the code.
They added multiple dungeons, a new display, special
individual character quests, a new endgame and many other new features, and
produced \textit{NetHack} 3.1.
Version 3.1.0 was released in January of 1993.

%.pg
\medskip
\noindent \textit{Ken Lorber}, \textit{Gregg Wonderly} and \textit{Greg Olson}, with help
from \textit{Richard Addison}, \textit{Mike Passaretti}, and \textit{Olaf Seibert},
developed \textit{NetHack} 3.1 for the Amiga.

%.pg
\medskip
\noindent \textit{Norm Meluch} and \textit{Kevin Smolkowski}, with help from
\textit{Carl Schelin}, \textit{Stephen Spackman}, \textit{Steve VanDevender},
and \textit{Paul Winner}, ported \textit{NetHack} 3.1 to the PC.

%.pg
\medskip
\noindent \textit{Jon W\textbraceleft tte} and \textit{Hao-yang Wang},
with help from \textit{Ross Brown}, \textit{Mike Engber}, \textit{David Hairston},
\textit{Michael Hamel}, \textit{Jonathan Handler}, \textit{Johnny Lee},
\textit{Tim Lennan}, \textit{Rob Menke}, and \textit{Andy Swanson},
developed \textit{NetHack} 3.1 for the Macintosh, porting it for MPW.
Building on their development, \textit{Bart House} added a Think C port.

%.pg
\medskip
\noindent \textit{Timo Hakulinen} ported \textit{NetHack} 3.1 to OS/2.
\textit{Eric Smith} ported \textit{NetHack} 3.1 to the Atari.
\textit{Pat Rankin}, with help from \textit{Joshua Delahunty},
was responsible for the VMS version of \textit{NetHack} 3.1.
\textit{Michael Allison} ported \textit{NetHack} 3.1 to Windows NT.

%.pg
\medskip
\noindent \textit{Dean Luick}, with help from \textit{David Cohrs}, developed
\textit{NetHack} 3.1 for X11.
It drew the map as text rather than graphically but
included \texttt{nh10.bdf}, an optionally used custom X11 font which has
tiny images in place of letters and punctuation, a precursor of tiles.
Those images don't extend to individual monster and object types, just
replacements for monster and object classes (so one custom image for all
``\texttt{a}'' insects and another for all ``\texttt{[}'' armor and so
forth, not separate images for beetles and ants or for cloaks and boots).

%.pg
\medskip
\noindent \textit{Warwick Allison} wrote a graphically displayed version
of \textit{NetHack}
for the Atari where the tiny pictures were described as ``icons'' and
were distinct for specific types of monsters and objects rather than just
their classes.
He contributed them to the \textit{NetHack Development Team} which
rechristened them ``tiles'', original usage which has subsequently been
picked up by various other games.
\textit{NetHack's} tiles support was then implemented on other platforms
(initially MS-DOS but eventually Windows, Qt, and X11 too).

%.pg
\medskip
\noindent The 3.2 \textit{NetHack Development Team}, comprised of \textit{Michael Allison}, \textit{Ken
Arromdee}, \textit{David Cohrs}, \textit{Jessie Collet}, \textit{Steve Creps}, {\it
Kevin Darcy}, \textit{Timo Hakulinen}, \textit{Steve Linhart}, \textit{Dean Luick},
\textit{Pat Rankin}, \textit{Eric Smith}, \textit{Mike Stephenson}, \textit{Janet Walz},
and \textit{Paul Winner}, released version 3.2.0 in April of 1996.

%.pg
\medskip
\noindent Version 3.2 marked the tenth anniversary of the formation of the
development team.
In a testament to their dedication to the game, all thirteen members
of the original \textit{NetHack Development Team} remained on the team at the
start of work on that release.
During the interval between the release of 3.1.3 and 3.2.0,
one of the founding members of the \textit{NetHack Development Team},
\textit{Dr. Izchak Miller}, was diagnosed with cancer and passed away.
That release of the game was
dedicated to him by the development and porting teams.

%.pg
Version 3.2 proved to be more stable than previous versions.
Many bugs were fixed, abuses eliminated, and game features tuned for
better game play.

%.pg
\medskip
During the lifespan of \textit{NetHack} 3.1 and 3.2, several enthusiasts
of the game added
their own modifications to the game and made these ``variants'' publicly
available:

%.pg
\medskip
\textit{Tom Proudfoot} and \textit{Yuval Oren} created \textit{NetHack++},
which was quickly renamed \textit{NetHack$--$}
when some people incorrectly assumed that it was a conversion of the
\textit{C} source code to \textit{C++}.
Working independently, \textit{Stephen White} wrote \textit{NetHack Plus}.
\textit{Tom Proudfoot} later merged \textit{NetHack Plus}
and his own \textit{NetHack$--$} to produce \textit{SLASH}.
\textit{Larry Stewart-Zerba} and \textit{Warwick Allison} improved the spell
casting system with the Wizard Patch.
\textit{Warwick Allison} also ported \textit{NetHack} to use the Qt interface.

%.pg
\medskip
\textit{Warren Cheung} combined \textit{SLASH} with the Wizard Patch
to produce \textit{Slash'EM}, and
with the help of \textit{Kevin Hugo}, added more features.
Kevin later joined the \textit{NetHack Development Team} and incorporated
the best of these ideas into \textit{NetHack} 3.3.

%.pg
\medskip
The final update to 3.2 was the bug fix release 3.2.3, which was released
simultaneously with 3.3.0 in December 1999 just in time for the Year 2000.
Because of the newer version, 3.2.3 was released as a source code patch only,
without any ready-to-play distribution for systems that usually had such.

%.pg
(To anyone considering resurrecting an old version:  all versions before
3.2.3 had a \textit{Y2K} bug.
The high scores file and the log file contained
dates which were formatted using a two-digit year, and 1999's year 99 was
followed by 2000's year 100.
That got written out successfully but it
unintentionally introduced an extra column in the file layout which prevented
score entries from being read back in correctly, interfering with insertion
of new high scores and with retrieval of old character names to use for
random ghost and statue names in the current game.)

%.pg
\medskip
The 3.3 \textit{NetHack Development Team}, consisting of \textit{Michael Allison}, \textit{Ken Arromdee},
\textit{David Cohrs}, \textit{Jessie Collet}, \textit{Steve Creps}, \textit{Kevin Darcy},
\textit{Timo Hakulinen}, \textit{Kevin Hugo}, \textit{Steve Linhart}, \textit{Ken Lorber},
\textit{Dean Luick}, \textit{Pat Rankin}, \textit{Eric Smith}, \textit{Mike Stephenson},
\textit{Janet Walz}, and \textit{Paul Winner}, released 3.3.0 in
December 1999 and 3.3.1 in August of 2000.

%.pg
\medskip
Version 3.3 offered many firsts. It was the first version to separate race
and profession. The Elf class was removed in preference to an elf race,
and the races of dwarves, gnomes, and orcs made their first appearance in
the game alongside the familiar human race.  Monk and Ranger roles joined
Archeologists, Barbarians, Cavemen, Healers, Knights, Priests, Rogues, Samurai,
Tourists, Valkyries and of course, Wizards.  It was also the first version
to allow you to ride a steed, and was the first version to have a publicly
available web-site listing all the bugs that had been discovered.  Despite
that constantly growing bug list, 3.3 proved stable enough to last for
more than a year and a half.

%.pg
\medskip
The 3.4 \textit{NetHack Development Team} initially consisted of
\textit{Michael Allison}, \textit{Ken Arromdee},
\textit{David Cohrs}, \textit{Jessie Collet}, \textit{Kevin Hugo}, \textit{Ken Lorber},
\textit{Dean Luick}, \textit{Pat Rankin}, \textit{Mike Stephenson},
\textit{Janet Walz}, and \textit{Paul Winner}, with \textit{ Warwick Allison} joining
just before the release of \textit{NetHack} 3.4.0 in March 2002.

%.pg
\medskip
As with version 3.3, various people contributed to the game as a whole as
well as supporting ports on the different platforms that \textit{NetHack}
runs on:

%.pg
\medskip
\noindent\textit{Pat Rankin} maintained 3.4 for VMS.

%.pg
\medskip
\noindent \textit{Michael Allison} maintained \textit{NetHack} 3.4 for the MS-DOS
platform.
\textit{Paul Winner} and \textit{Yitzhak Sapir} provided encouragement.

%.pg
\medskip
\noindent \textit{Dean Luick}, \textit{Mark Modrall}, and \textit{Kevin Hugo} maintained and
enhanced the Macintosh port of 3.4.

%.pg
\medskip
\noindent \textit{Michael Allison}, \textit{David Cohrs}, \textit{Alex Kompel},
\textit{Dion Nicolaas}, and
\textit{Yitzhak Sapir} maintained and enhanced 3.4 for the Microsoft Windows
platform.
\textit{Alex Kompel} contributed a new graphical interface for the Windows port.
\textit{Alex Kompel} also contributed a Windows CE port for 3.4.1.

%.pg
\medskip
\noindent \textit{Ron Van Iwaarden} was the sole maintainer of \textit{NetHack} for
OS/2 the past
several releases. Unfortunately Ron's last OS/2 machine stopped working in
early 2006. A great many thanks to Ron for keeping \textit{NetHack} alive on
OS/2 all these years.

%.pg
\medskip
\noindent \textit{Janne Salmij\"{a}rvi} and \textit{Teemu Suikki} maintained
and enhanced the Amiga port of 3.4 after \textit{Janne Salmij\"{a}rvi} resurrected
it for 3.3.1.

%.pg
\medskip
\noindent \textit{Christian ``Marvin'' Bressler} maintained 3.4 for the Atari after he
resurrected it for 3.3.1.

%.pg
\medskip
The release of \textit{NetHack} 3.4.3 in December 2003 marked the beginning of
a long release hiatus. 3.4.3 proved to be a remarkably stable version that
provided continued enjoyment by the community for more than a decade. The
\textit{NetHack Development Team} slowly and quietly continued to work on the game behind the scenes
during the tenure of 3.4.3. It was during that same period that several new
variants emerged within the \textit{NetHack} community. Notably sporkhack by
Derek S. Ray, \textit{unnethack} by Patric Mueller, \textit{nitrohack} and its
successors originally by Daniel Thaler and then by Alex Smith, and
\textit{Dynahack} by Tung Nguyen.
Some of those variants continue to be
developed, maintained, and enjoyed by the community to this day.

%.pg
\medskip
In September 2014, an interim snapshot of the code under development was
released publicly by other parties.
Since that code was a work-in-progress
and had not gone through the process of debugging it as a suitable release,
it was decided that the version numbers present on that code snapshot would
be retired and never used in an official \textit{NetHack} release.
An announcement was posted on the \textit{NetHack Development Team}'s official
\textit{nethack.org} website
to that effect, stating that there would never be a 3.4.4, 3.5, or 3.5.0
official release version.

%.pg
\medskip
In January 2015, preparation began for the release of \textit{NetHack} 3.6.

%.pg
\medskip
At the beginning of development for what would eventually get released as
3.6.0, the \textit{NetHack Development Team} consisted of \textit{Warwick Allison},
\textit{Michael Allison}, \textit{Ken Arromdee},
\textit{David Cohrs}, \textit{Jessie Collet},
\textit{Ken Lorber}, \textit{Dean Luick}, \textit{Pat Rankin},
\textit{Mike Stephenson}, \textit{Janet Walz}, and \textit{Paul Winner}.
In early 2015, ahead of the release of 3.6.0, new members
\textit{Sean Hunt}, \textit{Pasi Kallinen}, and \textit{Derek S. Ray}
joined the \textit{NetHack Development Team}.

%.pg
\medskip
Near the end of the development of 3.6.0, one of the significant inspirations
for many of the humorous and fun features found in the game,
author Terry Pratchett, passed away. \textit{NetHack} 3.6.0 introduced
a tribute to him.

%.pg
\medskip
3.6.0 was released in December 2015, and merged work done by the development
team since the release of 3.4.3 with some of the beloved community
patches.  Many bugs were fixed and some code was restructured.

%.pg
\medskip
The \textit{NetHack Development Team}, as well as \textit{Steve VanDevender} and
\textit{Kevin Smolkowski}, ensured that \textit{NetHack} 3.6 continued to
operate on various UNIX flavors and maintained the X11 interface.

%.pg
\medskip
\textit{Ken Lorber}, \textit{Haoyang Wang}, \textit{Pat Rankin}, and \textit{Dean Luick}
maintained the port of \textit{NetHack} 3.6 for MacOS.

%.pg
\medskip
\textit{Michael Allison}, \textit{David Cohrs}, \textit{Bart House},
\textit{Pasi Kallinen}, \textit{Alex Kompel}, \textit{Dion Nicolaas},
\textit{Derek S. Ray} and  \textit{Yitzhak Sapir}
maintained the port of  \textit{NetHack} 3.6 for Microsoft Windows.

%.pg
\medskip
\textit{Pat Rankin} attempted to keep the VMS port running for \textit{NetHack} 3.6,
hindered by limited access.  \textit{Kevin Smolkowski} has updated and tested it
for the most recent version of OpenVMS (V8.4 as of this writing) on Alpha
and Integrity (aka Itanium aka IA64) but not VAX.

%.pg
\medskip
\textit{Ray Chason}  resurrected the MS-DOS port for 3.6 and contributed the
necessary updates to the community at large.

%.pg
\medskip
In late April 2018, several hundred bug fixes for 3.6.0 and some new features
were assembled and released as \textit{NetHack} 3.6.1.
The \textit{NetHack Development Team} at the
time of release of 3.6.1 consisted of
\textit{Warwick Allison}, \textit{Michael Allison}, \textit{Ken Arromdee},
\textit{David Cohrs}, \textit{Jessie Collet},
\textit{Pasi Kallinen}, \textit{Ken Lorber}, \textit{Dean Luick},
\textit{Patric Mueller}, \textit{Pat Rankin}, \textit{Derek S. Ray},
\textit{Alex Smith}, \textit{Mike Stephenson}, \textit{Janet Walz}, and
\textit{Paul Winner}.

%.pg
\medskip
In early May 2019, another 320 bug fixes along with some enhancements and
the adopted curses window port, were released as 3.6.2.

%.pg
\medskip
\textit{Bart House}, who had contributed to the game as a porting team participant
for decades, joined the \textit{NetHack Development Team} in late May 2019.

%.pg
\medskip
\textit{NetHack} 3.6.3 was released on December 5, 2019 containing over 190 bug
fixes to \textit{NetHack} 3.6.2.

%.pg
\medskip
\textit{NetHack} 3.6.4 was released on December 18, 2019 containing a security fix and
a few bug fixes.

%.pg
\medskip
\textit{NetHack} 3.6.5 was released on January 27, 2020 containing some security fixes
and a small number of bug fixes.

%.pg
\medskip
\textit{NetHack} 3.6.6 was released on March 8, 2020 containing a security fix and
some bug fixes.

%.pg
\medskip
\textit{NetHack} 3.6.7 was released on February 16, 2023 containing a security fix and
some bug fixes.

%.pg
\medskip
Development work for the major release to follow \textit{NetHack} 3.6 began in 2015
around the same time as \textit{NetHack} 3.6.0 was being released. That development
work continued in parallel to each of the \textit{NetHack} 3.6 releases from 2015
through 2023, and continued until the end of April 2026.  For the first time,
that development was shared publicly on GitHub and SourceForge as it occurred.
It was done under the label NetHack-3.7 work-in-progress (WIP), although the
version number for the next release had not yet been solidified.

%.pg
\medskip
Exposure of the development to the public brought many good things, and some
challenges. People were able to observe and criticize changes and new features
almost immediately, and they often did. The GitHub pull request system made
it straightforward for people to contribute directly to development.
Contributions resolved many, many bugs in the game and we thank all the
contributors.

%.pg
\medskip
In early 2026, with the game development getting stable enough to consider
initiating an official release, the devteam reviewed the nature and number of
changes in the game. It was clear that there was sufficient depth and
breadth to warrant a major release and version 5.0 was decided on.
That's a new major release over 3.x, without opening up any ambiguity
or confusion with existing variants that there might have been had it
been released as version 4.0.

%.pg
\medskip
\textit{NetHack} 5.0.0 was released on May 2, 2026.

%.pg
\medskip
The source code for \textit{NetHack} 5.0.0 was modified and modernized to be
compliant with the C99 standard. The 5.0.0 release contained over 3100 fixes,
changes, and updated features.

%.pg
\medskip
\textit{NetHack} 5.0 was the first version to replace the lex and yacc level
and dungeon compilers of past versions, with a new Lua interpreter-based
approach to provide those elements. Lua is also used for the quest texts
in \textit{NetHack} 5.0.0. The entire development team acknowledges the work
done by \textit{Pasi Kallinen} to make that happen.

%.pg
\medskip
At the time of the \textit{NetHack} 5.0 release, the core
\textit{NetHack Development Team},
active and erstwhile, included \textit{Warwick Allison}, \textit{Michael Allison},
\textit{Ken Arromdee}, \textit{David Cohrs}, \textit{Jessie Collet}, \textit{Bart House},
\textit{Kevin Hugo}, \textit{Pasi Kallinen}, \textit{Ken Lorber}, \textit{Dean Luick},
\textit{Pat Rankin}, \textit{Derek S. Ray}, \textit{Alex Smith}, \textit{Patric Mueller},
\textit{Mike Stephenson}, \textit{Janet Walz}, \textit{Paul Winner}.

%.pg
\medskip
\textit{Ken Lorber}, \textit{Pat Rankin}, \textit{Patrick Mueller} and
\textit{Michael Allison} helped ensure that \textit{NetHack} 5.0 would run
on macOS.

%.pg
\medskip
\textit{Ingo Paschke} somehow managed to revive a \textit{NetHack} 5.0 port for the Amiga,
using a cross-compiler on a modern platform to do so. His work was shared so
others can straightforwardly produce \textit{NetHack} 5.0 for the Amiga thanks to his
efforts.

%.pg
\medskip
\textit{Ray Chason} contributed the majority of maintenance work for the
\textit{NetHack} 5.0 MS-DOS port, including porting the curses interface to it.
\textit{Michael Allison} ensured that \textit{NetHack} 5.0 core changes continued to
work with the msdos port and keep it alive.  Cross-compiling the MS-DOS
port has helped make that possible and mostly painless.

%.pg
\medskip
People that contributed to the Windows port of \textit{NetHack} 5.0 since the
development of \textit{NetHack} 5.0 began over eleven years ago, included
\textit{Michael Allison}, \textit{David Cohrs}, \textit{Bart House},
\textit{Pasi Kallinen}, \textit{Alex Kompel}, \textit{Dion Nicolaas},
\textit{Derek S. Ray} and \textit{Yitzhak Sapir}.

%.pg
\medskip
With sadness, the devteam would like to acknowledge and remember the past
contributions from the late \textit{Ron Van Iwaarden}, who was the sole maintainer
of \textit{NetHack} for OS/2 for several past \textit{NetHack} releases. Ron will be missed.

%.pg
\medskip
\noindent The official \textit{NetHack} web site is maintained by \textit{Ken Lorber} at
\url{https://www.nethack.org}.

%.pg
%.hn 2

\subsection*{Special Thanks}
\noindent On behalf of the \textit{NetHack} community, thank you very much once
again to \textit{M. Drew Streib} and \textit{Pasi Kallinen} for providing a
public \textit{NetHack} server at nethack.alt.org. Thanks to \textit{Keith Simpson}
and \textit{Andy Thomson} for hardfought.org. Thanks to all those
unnamed dungeoneers who invest their time and effort into annual
\textit{NetHack} tournaments such as \textit{Junethack},
\textit{The November NetHack Tournament}, and in days past,
\textit{devnull.net} (gone for now, but not forgotten).
\clearpage

%.hn 2
\subsection*{Dungeoneers}
%.pg
\noindent From time to time, some depraved individual out there in netland sends a
particularly intriguing modification to help out with the game.  The
\textit{NetHack Development Team} sometimes makes note of the names of the worst
of these miscreants in this, the list of Dungeoneers:
\medskip
%.sd
\begin{center}
\begin{tabular}{llll}
%TABLE_START
Adam Aronow & Irina Rempt-Drijfhout & Mike Gallop\\
Alex Kompel & Izchak Miller & Mike Passaretti\\
Alex Smith & J. Ali Harlow & Mike Stephenson\\
Andreas Dorn & Janet Walz & Mikko Juola\\
Andy Church & Janne Salmij\"{a}rvi & Nathan Eady\\
Andy Swanson & Jean-Christophe Collet & Norm Meluch\\
Andy Thomson & Jeff Bailey & Olaf Seibert\\
Ari Huttunen & Jochen Erwied & Pasi Kallinen\\
Bart House & John Kallen & Pat Rankin\\
Benson I. Margulies & John Rupley & Patric Mueller\\
Bill Dyer & John S. Bien & Paul Winner\\
Boudewijn Waijers & Johnny Lee & Pierre Martineau\\
Bruce Cox & Jon W\{tte & Ralf Brown\\
Bruce Holloway & Jonathan Handler & Ray Chason\\
Bruce Mewborne & Joshua Delahunty & Richard Addison\\
Cameron Root & Karl Garrison & Richard Beigel\\
Carl Schelin & Keizo Yamamoto & Richard P. Hughey\\
Chris Russo & Keith Simpson & Rob Menke\\
David Cohrs & Ken Arnold & Robin Bandy\\
David Damerell & Ken Arromdee & Robin Johnson\\
David Gentzel & Ken Lorber & Roderick Schertler\\
David Hairston & Ken Washikita & Roland McGrath\\
Dean Luick & Kestrel Gregorich-Trevor & Ron Van Iwaarden\\
Del Lamb & Kevin Darcy & Ronnen Miller\\
Derek S. Ray & Kevin Hugo & Ross Brown\\
Deron Meranda & Kevin Sitze & Sascha Wostmann\\
Dion Nicolaas & Kevin Smolkowski & Scott Bigham\\
Dylan O'Donnell & Kevin Sweet & Scott R. Turner\\
Eric Backus & Lars Huttar & Sean Hunt\\
Eric Hendrickson & Leon Arnott & Stephen Spackman\\
Eric R. Smith & M. Drew Streib & Stefan Thielscher\\
Eric S. Raymond & Malcolm Ryan & Stephen White\\
Erik Andersen & Mark Gooderum & Steve Creps\\
Fredrik Ljungdahl & Mark Modrall & Steve Linhart\\
Frederick Roeber & Marvin Bressler & Steve VanDevender\\
G. Branden Robinson & Matthew Day & Teemu Suikki\\
Gil Neiger & Merlyn LeRoy & Tim Lennan\\
Greg Laskin & Michael Allison & Timo Hakulinen\\
Greg Olson & Michael Feir & Tom Almy\\
Gregg Wonderly & Michael Hamel & Tom West\\
Hao-yang Wang & Michael Meyer & Warren Cheung\\
Helge Hafting & Michael Sokolov & Warwick Allison\\
Ingo Paschke & Mike Engber & Yitzhak Sapir
%TABLE_END  Do not delete this line.
\end{tabular}
\end{center}
%.ed
\clearpage
%\vfill
%\begin{flushleft}
%\small
%Microsoft and MS-DOS are registered trademarks of Microsoft Corporation.\\
%%%Don't need next line if a UNIX macro automatically inserts footnotes.
%UNIX is a registered trademark of AT\&T.\\
%Lattice is a trademark of Lattice, Inc.\\
%Atari and 1040ST are trademarks of Atari, Inc.\\
%Amiga is a trademark of Commodore-Amiga, Inc.\\
%%.sm
%Brand and product names are trademarks or registered trademarks
%of their respective holders.
%\end{flushleft}

\end{document}
